% ======================================================================
% Chapter 3: The Dirac Field
% (book pages 35-76; PDF pages 58-99)
% ======================================================================
\chapter{The Dirac Field}

Having exhaustively treated the simplest relativistic field equation, we now
move on to the second simplest, the Dirac equation.  You may already be
familiar with the Dirac equation in its original incarnation, that is, as a
single-particle quantum-mechanical wave equation.%
\footnote{This subject is covered, for example, in Schiff (1968), Chapter~13;
Baym (1969), Chapter~23; Sakurai (1967), Chapter~3.  Although the present
chapter is self-contained, we recommend that you also study the
single-particle Dirac equation at some point.}  In this chapter our viewpoint
will be quite different.  First we will rederive the Dirac equation as a
classical relativistic field equation, with special emphasis on its
relativistic invariance.  Then, in Section~3.5, we will quantize the Dirac
field in a manner similar to that used for the Klein--Gordon field.

\section{Lorentz Invariance in Wave Equations}

First we must address a question that we swept over in Chapter~2: What do we
mean when we say that an equation is ``relativistically invariant''?  A
reasonable definition is the following: If $\phi$ is a field or collection of
fields and $D$ is some differential operator, then the statement ``$D\phi =
0$ is relativistically invariant'' means that if $\phi(x)$ satisfies this
equation, and we perform a rotation or boost to a different frame of
reference, then the transformed field, in the new frame of reference,
satisfies the same equation.  Equivalently, we can imagine physically
rotating or boosting all particles or fields by a common angle or velocity;
again, the equation $D\phi = 0$ should be true after the transformation.  We
will adopt this ``active'' point of view toward transformations in the
following analysis.

The Lagrangian formulation of field theory makes it especially easy to
discuss Lorentz invariance.  An equation of motion is automatically Lorentz
invariant by the above definition if it follows from a Lagrangian that is a
Lorentz scalar.  This is an immediate consequence of the principle of least
action: If boosts leave the Lagrangian unchanged, the boost of an extremum
in the action will be another extremum.

As an example, consider the Klein--Gordon theory.  We can write an arbitrary
Lorentz transformation as
\begin{equation}
  x^\mu \to x'^{\mu} = \Lambda^\mu_{\ \nu} x^\nu,
  \label{eq:3.1}
\end{equation}
for some $4\times4$ matrix $\Lambda$.  What happens to the Klein--Gordon
field $\phi(x)$ under this transformation?  Think of the field $\phi$ as
measuring the local value of some quantity that is distributed through
space.  If there is an accumulation of this quantity at $x=x_0$, $\phi(x)$
will have a maximum at $x_0$.  If we now transform the original distribution
by a boost, the new distribution will have a maximum at $x=\Lambda x_0$.
This is illustrated in Fig.~\ref{fig:3.1}(a).  The corresponding
transformation of the field is
\begin{equation}
  \phi(x) \to \phi'(x) = \phi(\Lambda^{-1}x).
  \label{eq:3.2}
\end{equation}
That is, the transformed field, evaluated at the boosted point, gives the
same value as the original field evaluated at the point before boosting.

\begin{figure}[t]
  \centering
  \includegraphics[width=0.7\textwidth]{fig3.1.png}
  \caption{When a rotation is performed on a vector field, it affects the
    orientation of the vector as well as the location of the region
    containing the configuration.}
  \label{fig:3.1}
\end{figure}

We should check that this transformation leaves the form of the
Klein--Gordon Lagrangian unchanged.  According to \eqref{eq:3.2}, the mass
term $\frac{1}{2}m^2\phi^2(x)$ is simply shifted to the point
$(\Lambda^{-1}x)$.  The transformation of $\partial_\mu\phi(x)$ is
\begin{equation}
  \partial_\mu\phi(x) \to
    \frac{\partial}{\partial x^\mu}\bigl(\phi(\Lambda^{-1}x)\bigr)
    = (\Lambda^{-1})^{\nu}_{\ \mu}(\partial_\nu\phi)(\Lambda^{-1}x).
  \label{eq:3.3}
\end{equation}
Since the metric tensor $g_{\mu\nu}$ is Lorentz invariant, the matrices
$\Lambda^{-1}$ obey the identity
\begin{equation}
  (\Lambda^{-1})^{\mu}_{\ \rho}(\Lambda^{-1})^{\nu}_{\ \sigma}
    g_{\mu\nu} = g_{\rho\sigma} .
  \label{eq:3.4}
\end{equation}
Using this relation, we can compute the transformation law of the kinetic
term of the Klein--Gordon Lagrangian:
\begin{align}
  (\partial_\mu\phi(x))^2 &\to
    g^{\mu\nu}(\partial_\mu\phi'(x))(\partial_\nu\phi'(x))
  = g^{\mu\nu}\bigl[(\Lambda^{-1})^{\rho}_{\ \mu}\bigr]
    \bigl[(\Lambda^{-1})^{\sigma}_{\ \nu}\bigr]
    (\partial_\rho\phi)(\partial_\sigma\phi)(\Lambda^{-1}x)\\
  &= g^{\rho\sigma}(\partial_\rho\phi)(\partial_\sigma\phi)(\Lambda^{-1}x).
\end{align}
Thus, the whole Lagrangian is simply transformed as a scalar:
\begin{equation}
  \mathcal{L}(x) \to \mathcal{L}(\Lambda^{-1}x).
  \label{eq:3.5}
\end{equation}
The action $S$, formed by integrating $\mathcal{L}$ over spacetime, is
Lorentz invariant.  A similar calculation shows that the equation of motion
is invariant:
\begin{equation}
  (\partial^2 + m^2)\phi'(x) = \bigl[
    (\Lambda^{-1})^{\mu}_{\ \lambda}\partial_\lambda
    (\Lambda^{-1})^{\nu}_{\ \sigma}\partial_\sigma + m^2 \bigr]
    \phi(\Lambda^{-1}x) = (g^{\lambda\sigma}\partial_\lambda\partial_\sigma
    + m^2)\phi(\Lambda^{-1}x) = 0 .
\end{equation}
The transformation law \eqref{eq:3.2} used for $\phi$ is the simplest
possible transformation law for a field.  It is the only possibility for a
field that has just one component.  But we know examples of
multiple-component fields that transform in more complicated ways.  The
most familiar case is that of a vector field, such as the 4-current density
$j^\mu(x)$ or the vector potential $A^\mu(x)$.  In this case, the quantity
that is distributed in spacetime also carries an orientation, which must be
rotated or boosted.  As shown in Fig.~\ref{fig:3.1}(b), the orientation must
be rotated forward as the point of evaluation of the field is changed:
\begin{equation}
  V^\mu(x) \to (\Lambda^{-1})^{\mu}_{\ \nu}V^\nu(\Lambda^{-1}x)
  \qquad\text{under Lorentz transformations}.
\end{equation}
Tensors of arbitrary rank can be built out of vectors by adding more
indices, with correspondingly more factors of $\Lambda$ in the
transformation law.  Using such vector and tensor fields we can write a
variety of Lorentz-invariant equations, for example, Maxwell's equations,
\begin{equation}
  \partial_\mu F^{\mu\nu} = 0
  \qquad\text{or}\qquad
  \partial^2 A^\mu - \partial_\nu\partial^\mu A^\nu = 0,
  \label{eq:3.6}
\end{equation}
which follow from the Lagrangian
\begin{equation}
  \mathcal{L}_{\mathrm{Maxwell}} =
    -\tfrac{1}{4}(\partial_\mu A_\nu - \partial_\nu A_\mu)^2 .
  \label{eq:3.7}
\end{equation}
In general, any equation in which each term has the same set of uncontracted
Lorentz indices will naturally be invariant under Lorentz transformations.
This method of tensor notation yields a large class of Lorentz-invariant
equations, but it turns out that there are still more.  How do we find them?
We could try to systematically find all possible transformation laws for a
field.  Then it would not be hard to write invariant Lagrangians.  For
simplicity, we will restrict our attention to linear transformations, so
that, if $\Phi_a$ is an $n$-component multiplet, the Lorentz transformation
law is given by an $n\times n$ matrix $M(\Lambda)$:
\begin{equation}
  \Phi_a(x) \to M_{ab}(\Lambda)\Phi_b(\Lambda^{-1}x).
  \label{eq:3.8}
\end{equation}
It can be shown that the most general nonlinear transformation laws can be
built from these linear transformations, so there is no advantage in
considering transformations more general than \eqref{eq:3.8}.  In the
following discussion, we will suppress the change in the field argument and
write the transformation \eqref{eq:3.8} in the form
\begin{equation}
  \Phi \to M(\Lambda)\Phi .
  \label{eq:3.9}
\end{equation}
What are the possible allowed forms for the matrices $M(\Lambda)$?  The basic
restriction on $M(\Lambda)$ is found by imagining two successive
transformations, $\Lambda$ and $\Lambda'$.  The net result must be a new
Lorentz transformation $\Lambda''$; that is, the Lorentz transformations form
a group.  This gives a consistency condition that must be satisfied by the
matrices $M(\Lambda)$: Under the sequence of two transformations,
\begin{equation}
  \Phi \to M(\Lambda')M(\Lambda)\Phi = M(\Lambda'')\Phi,
  \label{eq:3.10}
\end{equation}
for $\Lambda'' = \Lambda'\Lambda$.  Thus the correspondence between the
matrices $M$ and the transformations $\Lambda$ must be preserved under
multiplication.  In mathematical language, we say that the matrices $M$ must
form an $n$-dimensional representation of the Lorentz group.  So our
question now is rephrased in mathematical language: What are the
(finite-dimensional) matrix representations of the Lorentz group?

Before answering this question for the Lorentz group, let us consider a
simpler group, the rotation group in three dimensions.  This group has
representations of every dimensionality $n$, familiar in quantum mechanics
as the matrices that rotate the $n$-component wavefunctions of particles of
different spins.  The dimensionality is related to the spin quantum number
$s$ by $n = 2s+1$.  The most important nontrivial representation is the
two-dimensional representation, corresponding to spin $1/2$.  The matrices of
this representation are the $2\times2$ unitary matrices with determinant 1,
which can be expressed as
\begin{equation}
  U = e^{-i\theta^l\sigma^l/2},
  \label{eq:3.11}
\end{equation}
where $\theta^l$ are three arbitrary parameters and $\sigma^l$ are the Pauli
sigma matrices.

For any continuous group, the transformations that lie infinitesimally close
to the identity define a vector space, called the Lie algebra of the group.
The basis vectors for this vector space are called the generators of the Lie
algebra, or of the group.  For the rotation group, the generators are the
angular momentum operators $J^i$, which satisfy the commutation relations
\begin{equation}
  [J^i, J^j] = i\epsilon^{ijk}J^k .
  \label{eq:3.12}
\end{equation}
The finite rotation operations are formed by exponentiating these operators:
In quantum mechanics, the operator
\begin{equation}
  R = \exp\Bigl[-\frac{i}{\hbar}\bm{\theta}\cdot\mathbf{J}\Bigr]
  \label{eq:3.13}
\end{equation}
gives the rotation by an angle $|\bm{\theta}|$ about the axis
$\bm{\theta}$.  The commutation relations of the operators $J^i$ determine
the multiplication laws of these rotation operators.  Thus, a set of matrices
satisfying the commutation relations \eqref{eq:3.12} produces, through
exponentiation as in \eqref{eq:3.13}, a representation of the rotation group.
In the example given in the previous paragraph, the representation of the
angular momentum operators
\begin{equation}
  J^l = \frac{\sigma^l}{2}
  \label{eq:3.14}
\end{equation}
produces the representation of the rotation group given in Eq.~\eqref{eq:3.11}.
It is generally true that one can find matrix representations of a continuous
group by finding matrix representations of the generators of the group (which
must satisfy the proper commutation relations), then exponentiating these
infinitesimal transformations.

For our present problem, we need to know the commutation relations of the
generators of the group of Lorentz transformations.  For the rotation group,
one can work out the commutation relations by writing the generators as
differential operators; from the expression
\begin{equation}
  \mathbf{J} = \mathbf{x}\times\mathbf{p} = \mathbf{x}\times(-i\nabla),
  \label{eq:3.15}
\end{equation}
the angular momentum commutation relations \eqref{eq:3.12} follow
straightforwardly.  The use of the cross product in \eqref{eq:3.15} is
special to the case of three dimensions.  However, we can also write the
operators as an antisymmetric tensor,
\begin{equation}
  J^i = -\tfrac{1}{2}\epsilon^{ijk}J^{jk},
\end{equation}
so that $J^3 = J^{12}$ and so on.  The generalization to four-dimensional
Lorentz transformations is now quite natural:
\begin{equation}
  J^{\mu\nu} = i(x^\mu\partial^\nu - x^\nu\partial^\mu).
  \label{eq:3.16}
\end{equation}
We will soon see that these six operators generate the three boosts and
three rotations of the Lorentz group.

To determine the commutation rules of the Lorentz algebra, we can now simply
compute the commutators of the differential operators \eqref{eq:3.16}.  The
result is
\begin{equation}
  [J^{\mu\nu}, J^{\rho\sigma}] = i\bigl( g^{\nu\rho}J^{\mu\sigma}
    - g^{\mu\rho}J^{\nu\sigma} - g^{\nu\sigma}J^{\mu\rho}
    + g^{\mu\sigma}J^{\nu\rho} \bigr).
  \label{eq:3.17}
\end{equation}
Any matrices that are to represent this algebra must obey these same
commutation rules.

Just to see that we have this right, let us look at one particular
representation (which we will simply pull out of a hat).  Consider the
$4\times4$ matrices
\begin{equation}
  (\mathcal{J}^{\mu\nu})^{\alpha}_{\ \beta} =
    i\bigl( \delta^{\mu\alpha}\delta^{\nu}_{\ \beta}
      - \delta^{\nu\alpha}\delta^{\mu}_{\ \beta} \bigr).
  \label{eq:3.18}
\end{equation}
It is not hard to check that these matrices satisfy the commutation relations
\eqref{eq:3.17}.  We can think of them as a set of linear transformations on
4-vectors, rotating or boosting a 4-vector $V^\mu$ in the way that we have
already seen.  A finite Lorentz transformation on a vector field is then
\begin{equation}
  V^\mu \to \Bigl[ \exp\bigl(
    -\tfrac{i}{2}\Omega_{\rho\sigma}\mathcal{J}^{\rho\sigma} \bigr)
    \Bigr]^{\mu}_{\ \nu} V^\nu ,
  \label{eq:3.19}
\end{equation}
where $\Omega_{\rho\sigma}$ are the parameters of the transformation, chosen
so that $\Lambda = \exp(-\frac{i}{2}\Omega_{\rho\sigma}\mathcal{J}^{\rho\sigma})$.

Now we need the corresponding representation for spinors.  The dimensionality
$n$ is 2 for spin $1/2$, so we seek $2\times2$ matrices that satisfy the same
commutation relations \eqref{eq:3.17}, with the generators representing
angular momentum $J^i = \sigma^i/2$ for the three-dimensional rotations.  We
can now see a systematic way to build the generators of the full Lorentz
algebra: Write the six generators in terms of the Pauli matrices, and then
look for a $2\times2$ representation.  The three rotation generators are
simply the spin operators
\begin{equation}
  S^{ij} = \frac{1}{2}\epsilon^{ijk}\sigma^k
  \label{eq:3.20}
\end{equation}
(compare Eq.~\eqref{eq:3.14}); this reproduces the commutators
\eqref{eq:3.12} for the rotation subgroup.  The three boost generators must
be $2\times2$ matrices; to find them, we note that the commutation relation
\begin{equation}
  [K^i, K^j] = -i\epsilon^{ijk}J^k
\end{equation}
(which follows from \eqref{eq:3.17}) must be satisfied.  One can check that
the matrices
\begin{equation}
  K^i = \frac{i}{2}\sigma^i
  \label{eq:3.21}
\end{equation}
satisfy this relation.  Thus a $2\times2$ representation of the Lorentz group
is
\begin{equation}
  \Lambda_{1/2} = \exp\Bigl( -\frac{i}{2}\Omega_{\mu\nu}S^{\mu\nu} \Bigr),
  \label{eq:3.22}
\end{equation}
where the generators $S^{\mu\nu}$ are given by
\begin{equation}
  S^{\mu\nu} = \begin{cases}
    \frac{1}{2}\epsilon^{ijk}\sigma^k & (\mu,\nu)=(i,j)\\
    \frac{i}{2}\sigma^i & (\mu,\nu)=(0,i)
  \end{cases}
\end{equation}
in terms of the Pauli matrices.  We note that $\sigma^i$ are Hermitian, so
the rotation generators $S^{ij}$ are Hermitian while the boost generators
$S^{0i}$ are anti-Hermitian.  The finite-dimensional representation
\eqref{eq:3.22} is therefore not unitary (as befits the noncompact Lorentz
group).  These spinor indices will not be our main concern, however; the
important point is that we now have a candidate $2\times2$ representation of
the Lorentz group that we can use to write the transformation law of a
2-component field.

Let us write out explicitly the objects that transform according to this
representation.  A Dirac spinor is a 4-component object
\begin{equation}
  \psi(x) = \begin{pmatrix}\psi_1\\ \psi_2\\ \psi_3\\ \psi_4\end{pmatrix},
\end{equation}
which we will often split into two 2-component pieces as we will describe
below.  Rather than building up the Lorentz algebra from scratch, we will use
a different, but equivalent, approach.  Define the four $4\times4$ matrices
$\gamma^\mu$ by the anticommutation relations
\begin{equation}
  \{\gamma^\mu, \gamma^\nu\} = 2g^{\mu\nu},
  \label{eq:3.22b}
\end{equation}
called the Dirac algebra.  (In Minkowski space the anticommutator is
$\{\gamma^\mu,\gamma^\nu\} = \gamma^\mu\gamma^\nu + \gamma^\nu\gamma^\mu$.)
One can show that the matrices
\begin{equation}
  S^{\mu\nu} = \frac{i}{4}[\gamma^\mu, \gamma^\nu]
  \label{eq:3.23}
\end{equation}
satisfy the commutation relations \eqref{eq:3.17} of the Lorentz algebra.
Thus the $\gamma$ matrices provide a representation of the Lorentz
generators, in which the transformation law of a Dirac field $\psi$ under a
Lorentz transformation is
\begin{equation}
  \psi(x) \to \Lambda_{1/2}\psi(\Lambda^{-1}x),
  \qquad \Lambda_{1/2} = \exp\Bigl(
    -\frac{i}{2}\Omega_{\mu\nu}S^{\mu\nu}\Bigr).
  \label{eq:3.24}
\end{equation}

For two-dimensional Minkowski space, we can find a set of $2\times2$
matrices satisfying the Dirac algebra:
\begin{equation}
  \gamma^0 = \sigma^3,\qquad \gamma^1 = i\sigma^1,
\end{equation}
so that
\begin{equation}
  \{\gamma^i, \gamma^j\} = -2\delta^{ij}.
  \label{eq:3.24b}
\end{equation}
The factor of $i$ in the first line and the minus sign in the second line are
purely conventional.  The matrices representing the Lorentz algebra are then
\begin{equation}
  S^{ij} = \frac{1}{2}\epsilon^{ijk}\sigma^k
  \label{eq:3.25}
\end{equation}
which we recognize as the two-dimensional representation of the rotation
group.

Now let us find Dirac matrices $\gamma^\mu$ for four-dimensional Minkowski
space.  It turns out that these matrices must be at least $4\times4$.  (There
is no fourth $2\times2$ matrix, for example, that anticommutes with the three
Pauli sigma matrices.)  Further, all $4\times4$ representations of the Dirac
algebra are unitarily equivalent.%
\footnote{This statement and the preceding one follow from the general theory
of the representations of the Lorentz group derived in Problem~3.1.}  We thus
need only write one explicit realization of the Dirac algebra.  One
representation, in $2\times2$ block form, is
\begin{equation}
  \gamma^0 = \begin{pmatrix}0&1\\1&0\end{pmatrix},\qquad
  \gamma^i = \begin{pmatrix}0&\sigma^i\\ -\sigma^i&0\end{pmatrix}.
  \label{eq:3.26}
\end{equation}
This representation is called the \emph{Weyl or chiral representation}.  We
will find it an especially convenient choice, and we will use it exclusively
throughout this book.  (Be careful, however, since many field theory
textbooks choose a different representation, in which $\gamma^0$ is
diagonal.  Furthermore, books that use chiral representations often make a
different choice of sign conventions.)

In our representation, the boost and rotation generators are
\begin{equation}
  S^{0i} = \frac{i}{2}\begin{pmatrix}-\sigma^i&0\\0&\sigma^i\end{pmatrix}
  \label{eq:3.27}
\end{equation}
and
\begin{equation}
  S^{ij} = \frac{1}{2}\epsilon^{ijk}
    \begin{pmatrix}\sigma^k&0\\0&\sigma^k\end{pmatrix}.
  \label{eq:3.28}
\end{equation}
A four-component field $\psi$ that transforms under boosts and rotations
according to \eqref{eq:3.27} and \eqref{eq:3.28} is called a \emph{Dirac
spinor}.  Note that the rotation generator $S^{ij}$ is just the
three-dimensional spinor transformation matrix \eqref{eq:3.25} replicated
twice.  The boost generators $S^{0i}$ are not Hermitian, and thus our
implementation of boosts is not unitary (this was also true of the vector
representation \eqref{eq:3.18}).  In fact the Lorentz group, being
``noncompact'', has no faithful, finite-dimensional representations that are
unitary.  But that does not matter to us, since $\psi$ is not a wavefunction;
it is a classical field.

Now that we have the transformation law for $\psi$, we should look for an
appropriate field equation.  One possibility is simply the Klein--Gordon
equation:
\begin{equation}
  (\partial^2 + m^2)\psi = 0 .
  \label{eq:3.29}
\end{equation}
This works because the spinor transformation matrices \eqref{eq:3.27} and
\eqref{eq:3.28} operate only in the ``internal'' space; they go right through
the differential operator.  But it is possible to write a stronger, first-order
equation, which implies \eqref{eq:3.29} but contains additional information.
To do this we need to know one more property of the $\gamma$ matrices.  With a
short computation you can verify that
\begin{equation}
  \{\gamma^\mu, S^{\nu\sigma}\} = (\mathcal{J}^{\nu\sigma})^{\mu}_{\ \rho}
    \gamma^\rho,
  \label{eq:3.30}
\end{equation}
or equivalently,
\begin{equation}
  \Bigl( 1 + \tfrac{i}{2}\omega_{\rho\sigma}S^{\rho\sigma} \Bigr)
    \gamma^\mu
    \Bigl( 1 - \tfrac{i}{2}\omega_{\rho\sigma}S^{\rho\sigma} \Bigr)
  = \Bigl( 1 - \tfrac{i}{2}\omega_{\rho\sigma}\mathcal{J}^{\rho\sigma}
    \Bigr)^{\mu}_{\ \nu}\gamma^\nu .
\end{equation}
This equation is just the infinitesimal form of
\begin{equation}
  \Lambda^{-1}_{1/2}\gamma^\mu\Lambda_{1/2}
    = \Lambda^{\mu}_{\ \nu}\gamma^\nu,
  \label{eq:3.31}
\end{equation}
where
\begin{equation}
  \Lambda_{1/2} = \exp\Bigl( -\frac{i}{2}\Omega_{\mu\nu}S^{\mu\nu} \Bigr)
  \label{eq:3.32}
\end{equation}
is the spinor representation of the Lorentz transformation $\Lambda$ (compare
\eqref{eq:3.19}).  Equation \eqref{eq:3.31} says that the $\gamma$ matrices
are invariant under simultaneous rotations of their vector and spinor indices
(just like the $\sigma^i$ under spatial rotations).  In other words, we can
``take the vector index $\mu$ on $\gamma^\mu$ seriously'', and dot
$\gamma^\mu$ into $\partial_\mu$ to form a Lorentz-invariant differential
operator.

We are now ready to write down the Dirac equation.  Here it is:
\begin{equation}
  (i\gamma^\mu\partial_\mu - m)\psi(x) = 0 .
  \label{eq:3.33}
\end{equation}
To show that it is Lorentz invariant, write down the Lorentz-transformed
version of the left-hand side and calculate:
\begin{align}
  (i\gamma^\mu\partial_\mu - m)\psi(x)
  &\to \bigl[ i\gamma^\mu(\Lambda^{-1})^{\nu}_{\ \mu}\partial_\nu - m \bigr]
    \Lambda_{1/2}\psi(\Lambda^{-1}x)\\
  &= \Lambda_{1/2}\bigl[ i\Lambda_{1/2}^{-1}\gamma^\mu\Lambda_{1/2}
    (\Lambda^{-1})^{\nu}_{\ \mu}\partial_\nu - m \bigr]\psi(\Lambda^{-1}x)\\
  &= \Lambda_{1/2}\bigl[ i\Lambda^{\mu}_{\ \rho}
    (\Lambda^{-1})^{\nu}_{\ \mu}\gamma^\rho\partial_\nu - m \bigr]
    \psi(\Lambda^{-1}x)\\
  &= \Lambda_{1/2}\bigl[ i\gamma^\rho\partial_\rho - m \bigr]
    \psi(\Lambda^{-1}x) = 0 .
\end{align}
To see that the Dirac equation implies the Klein--Gordon equation, act on the
left with $(-i\gamma^\nu\partial_\nu - m)$:
\begin{align}
  0 &= (-i\gamma^\nu\partial_\nu - m)(i\gamma^\mu\partial_\mu - m)\psi\\
  &= \Bigl( \gamma^\mu\gamma^\nu\partial_\mu\partial_\nu + m^2 \Bigr)\psi
    = \Bigl( \tfrac{1}{2}\{\gamma^\mu,\gamma^\nu\}\partial_\mu\partial_\nu
      + m^2 \Bigr)\psi
  = (\partial^2 + m^2)\psi .
\end{align}

To write down a Lagrangian for the Dirac theory, we must figure out how to
multiply two Dirac spinors to form a Lorentz scalar.  The obvious guess,
$\psi^\dagger\psi$, does not work.  Under a Lorentz boost, this becomes
$\psi^\dagger\Lambda_{1/2}^\dagger\Lambda_{1/2}\psi$; if the boost matrix
were unitary, we would have $\Lambda_{1/2}^\dagger = \Lambda_{1/2}^{-1}$ and
everything would be fine.  But $\Lambda_{1/2}$ is not unitary, because the
generators \eqref{eq:3.27} are not Hermitian.  The solution is to define
\begin{equation}
  \bar{\psi} = \psi^\dagger\gamma^0 .
  \label{eq:3.34}
\end{equation}
Under an infinitesimal Lorentz transformation parametrized by
$\Omega_{\mu\nu}$, we have
$\bar{\psi} \to \psi^\dagger\bigl( 1 + \frac{i}{2}\Omega_{\mu\nu}
(S^{\mu\nu})^\dagger \bigr)\gamma^0$.  The sum over $\mu$ and $\nu$ has six
distinct nonzero terms.  In the rotation terms, where $\mu$ and $\nu$ are
both nonzero, $(S^{\mu\nu})^\dagger = S^{\mu\nu}$ and commutes with
$\gamma^0$.  In the boost terms, where $\mu$ or $\nu$ is 0,
$(S^{0i})^\dagger = -S^{0i}$ but anticommutes with $\gamma^0$.  Passing the
$\gamma^0$ to the left therefore removes the dagger from $S^{\mu\nu}$, yielding
the transformation law
\begin{equation}
  \bar{\psi} \to \bar{\psi}\Lambda_{1/2}^{-1},
  \label{eq:3.35}
\end{equation}
and therefore the quantity $\bar{\psi}\psi$ is a Lorentz scalar.  Similarly
you can show (with the aid of \eqref{eq:3.31}) that $\bar{\psi}\gamma^\mu\psi$
is a Lorentz vector.

The correct, Lorentz-invariant Dirac Lagrangian is therefore
\begin{equation}
  \mathcal{L}_{\mathrm{Dirac}} = \bar{\psi}(i\gamma^\mu\partial_\mu
    - m)\psi .
  \label{eq:3.36}
\end{equation}
The Euler--Lagrange equation for $\bar{\psi}$ (or $\psi^\dagger$)
immediately yields the Dirac equation in the form \eqref{eq:3.33}; the
Euler--Lagrange equation for $\psi$ gives the same equation, in
Hermitian-conjugate form:
\begin{equation}
  i\partial_\mu\bar{\psi}\gamma^\mu + m\bar{\psi} = 0 .
  \label{eq:3.37}
\end{equation}

\subsection{Weyl Spinors}

From the block-diagonal form of the generators \eqref{eq:3.27} and
\eqref{eq:3.28}, it is apparent that the Dirac representation of the Lorentz
group is reducible.  We can form two 2-dimensional representations by
considering each block separately, and writing
\begin{equation}
  \psi = \begin{pmatrix}\psi_L\\ \psi_R\end{pmatrix}.
  \label{eq:3.38}
\end{equation}
The two-component objects $\psi_L$ and $\psi_R$ are called \emph{left-handed}
and \emph{right-handed Weyl spinors}.  You can easily verify that their
transformation laws, under infinitesimal rotations $\bm{\theta}$ and boosts
$\bm{\beta}$, are
\begin{align}
  \psi_L &\to \Bigl( 1 - i\bm{\theta}\cdot\tfrac{\bm{\sigma}}{2}
    - \bm{\beta}\cdot\tfrac{\bm{\sigma}}{2} \Bigr)\psi_L,\\
  \psi_R &\to \Bigl( 1 - i\bm{\theta}\cdot\tfrac{\bm{\sigma}}{2}
    + \bm{\beta}\cdot\tfrac{\bm{\sigma}}{2} \Bigr)\psi_R .
  \label{eq:3.39}
\end{align}
These transformation laws are connected by complex conjugation; using the
identity
\begin{equation}
  \sigma^2\bm{\sigma}^{*} = -\bm{\sigma}\sigma^2,
  \label{eq:3.40}
\end{equation}
it is not hard to show that the quantity $\sigma^2\psi_L^{*}$ transforms like
a right-handed spinor.

In terms of $\psi_L$ and $\psi_R$, the Dirac equation is
\begin{equation}
  \begin{pmatrix}
    -m & i\sigma^\mu\partial_\mu\\
    i\bar{\sigma}^\mu\partial_\mu & -m
  \end{pmatrix}
  \begin{pmatrix}\psi_L\\ \psi_R\end{pmatrix} = 0 .
  \label{eq:3.41}
\end{equation}
The two Lorentz group representations $\psi_L$ and $\psi_R$ are mixed by the
mass term in the Dirac equation.  But if we set $m=0$, the equations for
$\psi_L$ and $\psi_R$ decouple:
\begin{align}
  i(\partial_0 - \bm{\sigma}\cdot\nabla)\psi_L &= 0;\\
  i(\partial_0 + \bm{\sigma}\cdot\nabla)\psi_R &= 0 .
  \label{eq:3.42}
\end{align}
These are called the \emph{Weyl equations}; they are especially important
when treating neutrinos and the theory of weak interactions.

It is possible to clean up this notation slightly.  Define
\begin{equation}
  \sigma^\mu = (1, \bm{\sigma}),\qquad
  \bar{\sigma}^\mu = (1, -\bm{\sigma}),
  \label{eq:3.43}
\end{equation}
so that
\begin{equation}
  \gamma^\mu = \begin{pmatrix}
    0 & \sigma^\mu\\
    \bar{\sigma}^\mu & 0
  \end{pmatrix}.
  \label{eq:3.44}
\end{equation}
(The bar on $\sigma$ has absolutely nothing to do with the bar on $\psi$.)
Then the Dirac equation can be written
\begin{equation}
  \begin{pmatrix}
    -m & i\sigma\cdot\partial\\
    i\bar{\sigma}\cdot\partial & -m
  \end{pmatrix}
  \begin{pmatrix}\psi_L\\ \psi_R\end{pmatrix} = 0
  \label{eq:3.45}
\end{equation}
and the Weyl equations become
\begin{equation}
  i\bar{\sigma}\cdot\partial\,\psi_L = 0;\qquad
  i\sigma\cdot\partial\,\psi_R = 0 .
  \label{eq:3.46}
\end{equation}

\section{Free-Particle Solutions of the Dirac Equation}

To get some feel for the physics of the Dirac equation, let us now discuss
its plane-wave solutions.  Since a Dirac field $\psi$ obeys the
Klein--Gordon equation, we know immediately that it can be written as a
linear combination of plane waves:
\begin{equation}
  \psi(x) = u(p)e^{-ip\cdot x},\qquad \text{where } p^2 = m^2 .
  \label{eq:3.47}
\end{equation}
For the moment we will concentrate on solutions with positive frequency, that
is, $p^0>0$.  The column vector $u(p)$ must obey an additional constraint,
found by plugging \eqref{eq:3.47} into the Dirac equation:
\begin{equation}
  (i\gamma^\mu\partial_\mu - m)u(p)e^{-ip\cdot x} = 0
  \qquad\Rightarrow\qquad
  (\gamma^\mu p_\mu - m)u(p) = 0 .
  \label{eq:3.48}
\end{equation}
It is easiest to analyze this equation in the rest frame, where
$p = p_0 = (m, \mathbf{0})$; the solution for general $p$ can then be found by
boosting with $\Lambda_{1/2}$.  In the rest frame, Eq.~\eqref{eq:3.48}
becomes
\begin{equation}
  (m\gamma^0 - m)u(p_0) = m
    \begin{pmatrix}-1&1\\1&-1\end{pmatrix}u(p_0) = 0 ,
\end{equation}
and the solutions are
\begin{equation}
  u(p_0) = \sqrt{m}\begin{pmatrix}\xi\\ \xi\end{pmatrix},
  \label{eq:3.49}
\end{equation}
for any numerical two-component spinor $\xi$.  We conventionally normalize
$\xi$ so that $\xi^\dagger\xi = 1$; the factor $\sqrt{m}$ has been inserted
for future convenience.  We can interpret the spinor $\xi$ by looking at the
rotation generator \eqref{eq:3.28}: $\xi$ transforms under rotations as an
ordinary two-component spinor of the rotation group, and therefore determines
the spin orientation of the Dirac solution in the usual way.  For example,
when $\xi = \binom{1}{0}$, the particle has spin up along the 3-direction.

Notice that after applying the Dirac equation, we are free to choose only two
of the four components of $u(p)$.  This is just what we want, since a spin-$1/2$
particle has only two physical states---spin up and spin down.  (Of course we
are being a bit premature in talking about particles and spin.  We will prove
that the spin angular momentum of a Dirac particle is $\hbar/2$ when we
quantize the Dirac theory in Section~3.5; for now, just notice that there are
two possible solutions $u(p)$ for any momentum $p$.)

Now that we have the general form of $u(p)$ in the rest frame, we can obtain
$u(p)$ in any other frame by boosting.  Consider a boost along the 3-direction.
First we should remind ourselves of what the boost does to the 4-momentum
vector.  In infinitesimal form,
\begin{equation}
  p^\mu \to (1 - i\eta K^3)^{\mu}_{\ \nu}p^\nu
    = \begin{pmatrix}
        \cosh\eta & \sinh\eta\\
        \sinh\eta & \cosh\eta
      \end{pmatrix}
      \begin{pmatrix}m\\ 0\end{pmatrix}
    = \begin{pmatrix}m\cosh\eta\\ m\sinh\eta\end{pmatrix},
\end{equation}
where $\eta$ is some infinitesimal parameter.  For finite $\eta$ we must write
\begin{equation}
  p^\mu = \begin{pmatrix}
    \cosh\eta & \sinh\eta\\
    \sinh\eta & \cosh\eta
  \end{pmatrix}
  \begin{pmatrix}m\\ 0\end{pmatrix}
  = \begin{pmatrix}
    E\\ p^3
  \end{pmatrix},
  \label{eq:3.50}
\end{equation}
with $E = m\cosh\eta$, $p^3 = m\sinh\eta$.  The parameter $\eta$ is called the
\emph{rapidity}.  It is the quantity that is additive under successive
boosts.

Now apply the same boost to $u(p)$.  According to Eqs.~\eqref{eq:3.27} and
\eqref{eq:3.32},
\begin{align}
  u(p) &= \exp\Bigl( \tfrac{1}{2}\eta\gamma^0\gamma^3 \Bigr)
    \sqrt{m}\begin{pmatrix}\xi\\ \xi\end{pmatrix}
  = \begin{pmatrix}
      e^{\eta/2} & 0\\
      0 & e^{-\eta/2}
    \end{pmatrix}
    \sqrt{m}\begin{pmatrix}\xi\\ \xi\end{pmatrix}\\
  &= \sqrt{m}
    \begin{pmatrix}
      e^{\eta/2}\xi\\
      e^{-\eta/2}\xi
    \end{pmatrix}
  = \begin{pmatrix}
      \sqrt{p\cdot\sigma}\,\xi\\
      \sqrt{p\cdot\bar{\sigma}}\,\xi
    \end{pmatrix}.
  \label{eq:3.51}
\end{align}
The last line can be simplified to give
\begin{equation}
  u(p) = \begin{pmatrix}
    \sqrt{p\cdot\sigma}\,\xi\\
    \sqrt{p\cdot\bar{\sigma}}\,\xi
  \end{pmatrix},
  \label{eq:3.52}
\end{equation}
where it is understood that in taking the square root of a matrix, we take
the positive root of each eigenvalue.  This expression for $u(p)$ is not only
more compact, but is also valid for an arbitrary direction of $\mathbf{p}$.
When working with expressions of this form, it is often useful to know the
identity
\begin{equation}
  \sqrt{p\cdot\sigma}\,\sqrt{p\cdot\bar{\sigma}} = m
  \qquad\text{(and similarly with $\sigma$ and $\bar{\sigma}$ interchanged)}.
  \label{eq:3.53}
\end{equation}

For $\xi = \binom{1}{0}$ (spin up along the 3-axis), we have
\begin{equation}
  u(p) = \begin{pmatrix}
    \sqrt{p^0-p^3}\,\binom{1}{0}\\
    \sqrt{p^0+p^3}\,\binom{1}{0}
  \end{pmatrix}
  = \begin{pmatrix}
    \sqrt{E-p^3}\binom{1}{0}\\
    \sqrt{E+p^3}\binom{1}{0}
  \end{pmatrix},
  \label{eq:3.54}
\end{equation}
while for $\xi = \binom{0}{1}$ (spin down along the 3-axis) we have
\begin{equation}
  u(p) = \begin{pmatrix}
    \sqrt{p\cdot\sigma}\,\binom{0}{1}\\
    \sqrt{p\cdot\bar{\sigma}}\,\binom{0}{1}
  \end{pmatrix}
  = \begin{pmatrix}
    \sqrt{E+p^3}\binom{0}{1}\\
    \sqrt{E-p^3}\binom{0}{1}
  \end{pmatrix}.
  \label{eq:3.55}
\end{equation}
In the limit $\eta\to\infty$ the states degenerate into the two-component
spinors of a massless particle.  (We now see the reason for the factor of
$\sqrt{m}$ in \eqref{eq:3.49}: It keeps the spinor expressions finite in the
massless limit.)

The solutions \eqref{eq:3.54} and \eqref{eq:3.55} are eigenstates of the
helicity operator,
\begin{equation}
  h = \frac{1}{2}\hat{\mathbf{p}}\cdot\bm{\sigma}
  = \frac{1}{2}\hat{\mathbf{p}}\cdot\mathbf{J}.
  \label{eq:3.56}
\end{equation}
A particle with $h = +1/2$ is called right-handed, while one with $h =
-1/2$ is called left-handed.  The helicity of a massive particle depends on
the frame of reference, since one can always boost to a frame in which its
momentum is in the opposite direction (but its spin is unchanged).  For a
massless particle, which travels at the speed of light, one cannot perform
such a boost.

The extremely simple form of $u(p)$ for a massless particle in a helicity
eigenstate makes the behavior of such a particle easy to understand.  In
Chapter~1, it enabled us to guess the form of the
$e^{+}e^{-}\to\mu^{+}\mu^{-}$ cross section in the massless limit.  In
subsequent chapters we will often do a mindless calculation first, then look
at helicity eigenstates in the high-energy limit to understand what we have
done.

Incidentally, we are now ready to understand the origin of the notation
$\psi_L$ and $\psi_R$ for Weyl spinors.  The solutions of the Weyl equations
are states of definite helicity, corresponding to left- and right-handed
particles, respectively.  The Lorentz invariance of helicity (for a massless
particle) is manifest in the notation of Weyl spinors, since $\psi_L$ and
$\psi_R$ live in different representations of the Lorentz group.

It is convenient to write the normalization condition for $u(p)$ in a
Lorentz-invariant way.  We saw above that $u^\dagger u$ is not Lorentz
invariant.  Similarly,
\begin{equation}
  u^\dagger(p)u(p) = 2E_p\,\xi^\dagger\xi .
  \label{eq:3.57}
\end{equation}
To make a Lorentz scalar we define
\begin{equation}
  \bar{u}(p) = u^\dagger(p)\gamma^0 .
  \label{eq:3.58}
\end{equation}
Then by an almost identical calculation,
\begin{equation}
  \bar{u}u = 2m\,\xi^\dagger\xi .
  \label{eq:3.59}
\end{equation}
This will be our normalization condition, once we also require that the
two-component spinor $\xi$ be normalized as usual: $\xi^\dagger\xi = 1$.  It
is also conventional to choose basis spinors $\xi^1$ and $\xi^2$ (such as
$\binom{1}{0}$ and $\binom{0}{1}$) that are orthogonal.  For a massless
particle Eq.~\eqref{eq:3.59} is trivial, so we must write the normalization
condition in the form of \eqref{eq:3.57}.

Let us summarize our discussion so far.  The general solution of the Dirac
equation can be written as a linear combination of plane waves.  The
positive-frequency waves are of the form
\begin{equation}
  \psi(x) = u(p)e^{-ip\cdot x},\qquad p^2 = m^2,\quad p^0>0 .
  \label{eq:3.60}
\end{equation}
There are two linearly independent solutions for $u(p)$,
\begin{equation}
  u^s(p),\qquad s = 1, 2,
  \label{eq:3.61}
\end{equation}
which we normalize according to
\begin{equation}
  \bar{u}^r(p)u^s(p) = 2m\delta^{rs}
  \qquad\text{or}\qquad
  u^{r\dagger}(p)u^s(p) = 2E_p\delta^{rs}.
  \label{eq:3.62}
\end{equation}
In exactly the same way, we can find the negative-frequency solutions:
\begin{equation}
  \psi(x) = v(p)e^{+ip\cdot x},\qquad p^2 = m^2,\quad p^0>0 .
  \label{eq:3.63}
\end{equation}
(Note that we have chosen to put the $+$ sign into the exponential, rather
than having $p^0<0$.)  There are two linearly independent solutions for
$v(p)$,
\begin{equation}
  v^s(p),\qquad s = 1, 2,
  \label{eq:3.64}
\end{equation}
where $\eta^s$ is another basis of two-component spinors.  These solutions
are normalized according to
\begin{equation}
  \bar{v}^r(p)v^s(p) = -2m\delta^{rs}
  \qquad\text{or}\qquad
  v^{r\dagger}(p)v^s(p) = +2E_p\delta^{rs}.
  \label{eq:3.65}
\end{equation}
The $u$'s and $v$'s are also orthogonal to each other:
\begin{equation}
  \bar{u}^r(p)v^s(p) = \bar{v}^r(p)u^s(p) = 0 .
  \label{eq:3.66}
\end{equation}
Be careful, since $u^{r\dagger}(p)v^s(p)\neq 0$ and
$v^{r\dagger}(p)u^s(p)\neq 0$.  However, note that
\begin{equation}
  u^{r\dagger}(p)v^s(-\mathbf{p}) =
  v^{r\dagger}(-\mathbf{p})u^s(p) = 0 ,
  \label{eq:3.67}
\end{equation}
where we have changed the sign of the 3-momentum in one factor of each spinor
product.

\subsection{Spin Sums}

In evaluating Feynman diagrams, we will often wish to sum over the
polarization states of a fermion.  We can derive the relevant completeness
relations with a simple calculation:
\begin{align}
  \sum_{s=1,2} u^s(p)\bar{u}^s(p) &=
    \begin{pmatrix}
      \sqrt{p\cdot\sigma}\,\xi^s\\
      \sqrt{p\cdot\bar{\sigma}}\,\xi^s
    \end{pmatrix}
    \Bigl( (\xi^{s\dagger}\sqrt{p\cdot\bar{\sigma}})
      \quad (\xi^{s\dagger}\sqrt{p\cdot\sigma}) \Bigr)\\
  &= \begin{pmatrix}
      \sqrt{p\cdot\sigma}\,(\sum_s \xi^s\xi^{s\dagger})\sqrt{p\cdot\bar{\sigma}}
        & \sqrt{p\cdot\sigma}\,(\sum_s \xi^s\xi^{s\dagger})\sqrt{p\cdot\sigma}\\
      \sqrt{p\cdot\bar{\sigma}}\,(\sum_s \xi^s\xi^{s\dagger})
        \sqrt{p\cdot\bar{\sigma}}
        & \sqrt{p\cdot\bar{\sigma}}\,(\sum_s \xi^s\xi^{s\dagger})
          \sqrt{p\cdot\sigma}
    \end{pmatrix}\\
  &= \begin{pmatrix}
      \sqrt{p\cdot\sigma}\,\mathbf{1}\sqrt{p\cdot\bar{\sigma}}
        & \sqrt{p\cdot\sigma}\,\mathbf{1}\sqrt{p\cdot\sigma}\\
      \sqrt{p\cdot\bar{\sigma}}\,\mathbf{1}\sqrt{p\cdot\bar{\sigma}}
        & \sqrt{p\cdot\bar{\sigma}}\,\mathbf{1}\sqrt{p\cdot\sigma}
    \end{pmatrix}
  = \begin{pmatrix}
      m & p\cdot\sigma\\
      p\cdot\bar{\sigma} & m
    \end{pmatrix}.
  \label{eq:3.68}
\end{align}
In the second line we have used
\begin{equation}
  \sum_{s=1,2}\xi^s\xi^{s\dagger} = \begin{pmatrix}1&0\\0&1\end{pmatrix}.
\end{equation}
Thus we arrive at the desired formula,
\begin{equation}
  \sum_{s} u^s(p)\bar{u}^s(p) = \gamma\cdot p + m .
  \label{eq:3.69}
\end{equation}
Similarly,
\begin{equation}
  \sum_{s} v^s(p)\bar{v}^s(p) = \gamma\cdot p - m .
  \label{eq:3.70}
\end{equation}
The combination $\gamma\cdot p$ occurs so often that Feynman introduced the
notation
\begin{equation}
  \slashed{p} = \gamma^\mu p_\mu .
  \label{eq:3.71}
\end{equation}
We will use this notation frequently from now on.

\section{Dirac Matrices and Dirac Field Bilinears}

We saw in Section~3.2 that the quantity $\bar{\psi}\psi$ is a Lorentz scalar.
It is also easy to show that $\bar{\psi}\gamma^\mu\psi$ is a 4-vector---we used
this fact in writing down the Dirac Lagrangian \eqref{eq:3.36}.  Now let us
ask a more general question: Consider the expression $\bar{\psi}\Gamma\psi$,
where $\Gamma$ is any $4\times4$ constant matrix.  Can we decompose this
expression into terms that have definite transformation properties under the
Lorentz group?  The answer is yes, if we write $\Gamma$ in terms of the
following basis of sixteen $4\times4$ matrices, defined as antisymmetric
combinations of $\gamma$-matrices:
\begin{equation}
  \begin{array}{lll}
    1 & \text{(1 of these)} & \text{scalar}\\
    \gamma^\mu & \text{(4 of these)} & \text{vector}\\
    \sigma^{\mu\nu} = \tfrac{i}{2}[\gamma^\mu,\gamma^\nu] &
      \text{(6 of these)} & \text{tensor}\\
    \gamma^\mu\gamma^5 & \text{(4 of these)} & \text{pseudo-vector}\\
    \gamma^5 & \text{(1 of these)} & \text{pseudo-scalar}
  \end{array}
\end{equation}
The Lorentz-transformation properties of these matrices are easy to
determine.  For example,
\begin{align}
  \bar{\psi}[\gamma^\mu,\gamma^\nu]\psi
  &\to (\psi^\dagger\Lambda_{1/2}^\dagger)
    \tfrac{i}{2}[\gamma^\mu,\gamma^\nu](\Lambda_{1/2}\psi)\\
  &= \tfrac{i}{2}\Lambda^{\mu}_{\ \alpha}\Lambda^{\nu}_{\ \beta}
    \psi^\dagger[\gamma^\alpha,\gamma^\beta]\psi .
\end{align}
Each set of matrices transforms as an antisymmetric tensor of successively
higher rank.

The last two sets of matrices can be simplified by introducing an additional
gamma matrix,
\begin{equation}
  \gamma^5 = i\gamma^0\gamma^1\gamma^2\gamma^3
    = -\frac{i}{4!}\epsilon^{\mu\nu\rho\sigma}
      \gamma_\mu\gamma_\nu\gamma_\rho\gamma_\sigma .
  \label{eq:3.72}
\end{equation}
Then $\sigma^{\mu\nu}\gamma^5 = -i\epsilon^{\mu\nu\rho\sigma}\sigma_{\rho\sigma}$
and $\gamma^\mu\gamma^5 = +i\epsilon^{\mu\nu\rho\sigma}\sigma_{\nu\rho}\gamma_\sigma$.
The matrix $\gamma^5$ has the following properties, all of which can be
verified using \eqref{eq:3.72} and the anticommutation relations
\eqref{eq:3.22b}:
\begin{align}
  (\gamma^5)^\dagger &= \gamma^5; \label{eq:3.73}\\
  (\gamma^5)^2 &= 1; \label{eq:3.74}\\
  \{\gamma^5, \gamma^\mu\} &= 0 . \label{eq:3.75}
\end{align}
This last property implies that $[\gamma^5, S^{\mu\nu}] = 0$.  Thus the Dirac
representation must be reducible, since eigenvectors of $\gamma^5$ whose
eigenvalues are different transform without mixing (this criterion for
reducibility is known as Schur's lemma).  In our basis,
\begin{equation}
  \gamma^5 = \begin{pmatrix}-1&0\\0&1\end{pmatrix}
  \label{eq:3.76}
\end{equation}
in block-diagonal form.  So a Dirac spinor with only left- (right-) handed
components is an eigenstate of $\gamma^5$ with eigenvalue $-1$ ($+1$), and
indeed these spinors do transform without mixing, as we saw explicitly in
Section~3.2.

Let us now rewrite our table of $4\times4$ matrices, and introduce some
standard terminology:
\begin{equation}
  \begin{array}{lll}
    \text{scalar} & 1 & \text{(1)}\\
    \text{vector} & \gamma^\mu & \text{(4)}\\
    \text{tensor} & \sigma^{\mu\nu} = \tfrac{i}{2}[\gamma^\mu,\gamma^\nu]
      & \text{(6)}\\
    \text{pseudo-vector} & \gamma^\mu\gamma^5 & \text{(4)}\\
    \text{pseudo-scalar} & \gamma^5 & \text{(1)}
  \end{array}
  \label{eq:3.77}
\end{equation}
The terms pseudo-vector and pseudo-scalar arise from the fact that these
quantities transform as a vector and scalar, respectively, under continuous
Lorentz transformations, but with an additional sign change under parity
transformations (as we will discuss in Section~3.6).

From the vector and pseudo-vector matrices we can form two currents out of
Dirac field bilinears:
\begin{equation}
  j^\mu(x) = \bar{\psi}(x)\gamma^\mu\psi(x);\qquad
  j^{\mu 5}(x) = \bar{\psi}(x)\gamma^\mu\gamma^5\psi(x) .
  \label{eq:3.78}
\end{equation}
Let us compute the divergences of these currents, assuming that $\psi$
satisfies the Dirac equation:
\begin{align}
  \partial_\mu j^\mu &= (\partial_\mu\bar{\psi})\gamma^\mu\psi
    + \bar{\psi}\gamma^\mu(\partial_\mu\psi)\\
  &= (im\bar{\psi})\psi + \bar{\psi}(-im\psi) = 0 .
  \label{eq:3.79}
\end{align}
Thus $j^\mu$ is always conserved if $\psi(x)$ satisfies the Dirac equation.
When we couple the Dirac field to the electromagnetic field, $j^\mu$ will
become the electric current density.  Similarly, one can compute
\begin{equation}
  \partial_\mu j^{\mu 5} = 2im\bar{\psi}\gamma^5\psi .
  \label{eq:3.80}
\end{equation}
If $m=0$, this current (often called the axial vector current) is also
conserved.  It is then useful to form the linear combinations
\begin{equation}
  j_L^\mu = \bar{\psi}\gamma^\mu\frac{1-\gamma^5}{2}\psi
    = \bar{\psi}_L\gamma^\mu\psi_L;\qquad
  j_R^\mu = \bar{\psi}\gamma^\mu\frac{1+\gamma^5}{2}\psi
    = \bar{\psi}_R\gamma^\mu\psi_R .
  \label{eq:3.81}
\end{equation}
When $m=0$, these are the electric current densities of left-handed and
right-handed particles, respectively, and are separately conserved.

The two currents $j^\mu(x)$ and $j^{\mu 5}(x)$ are the Noether currents
corresponding to the two transformations
\begin{equation}
  \psi(x) \to e^{i\alpha}\psi(x)
  \qquad\text{and}\qquad
  \psi(x) \to e^{i\alpha\gamma^5}\psi(x) .
\end{equation}
The first of these is a symmetry of the Dirac Lagrangian \eqref{eq:3.36}.
The second, called a \emph{chiral transformation}, is a symmetry of the
derivative term in $\mathcal{L}$ but not the mass term; thus, Noether's
theorem confirms that the axial vector current is conserved only if $m=0$.

Products of Dirac bilinears obey interchange relations, known as Fierz
identities.  We will discuss only the simplest of these, which will be needed
several times later in the book.  This simplest identity is most easily
written in terms of the two-component Weyl spinors introduced in
Eq.~\eqref{eq:3.38}.  The core of the relation is the identity for the
$2\times2$ matrices $\sigma^\mu$ defined in Eq.~\eqref{eq:3.43}:
\begin{equation}
  (\sigma^\mu)_{\alpha\dot{\alpha}}(\sigma_\mu)_{\beta\dot{\beta}}
    = 2\epsilon_{\alpha\beta}\epsilon_{\dot{\alpha}\dot{\beta}} .
  \label{eq:3.82}
\end{equation}
(Here $\alpha, \beta,$ etc.\ are spinor indices, and $\epsilon$ is the
antisymmetric symbol.)  One can understand this relation by noting that the
indices $\alpha, \gamma$ transform in the Lorentz representation of $\psi_L$,
while $\beta, \delta$ transform in the separate representation of $\psi_R$,
and the whole quantity must be a Lorentz invariant.  Alternatively, one can
just verify the 16 components of \eqref{eq:3.82} explicitly.

By sandwiching identity \eqref{eq:3.82} between the right-handed portions
(i.e., lower half) of Dirac spinors $u_1, u_2, u_3, u_4$, we find the identity
\begin{equation}
  (\bar{u}_{1R}\sigma^\mu u_{2R})(\bar{u}_{3R}\sigma_\mu u_{4R})
    = 2\epsilon_{\alpha\beta}\,
    (\bar{u}_{1R})_\alpha(\bar{u}_{3R})_\beta\,
    \epsilon_{\gamma\delta}\,
    (u_{2R})_\gamma(u_{4R})_\delta ,
  \label{eq:3.83}
\end{equation}
which is one of the standard Fierz identities.

\section{Quantization of the Dirac Field}

We are now ready to quantize the Dirac field, following the same procedure we
used for the Klein--Gordon field in Section~2.3.  We will promote the
classical field $\psi$ and its conjugate momentum to operators obeying
canonical commutation relations.  In the case of the Dirac field, however, we
must impose \emph{anticommutation} relations in order that the theory have
only positive-energy excitations.  The formal justification of this
procedure---the spin-statistics theorem---will be discussed below, following
the same logic that led us to Fermi statistics in the previous section.

The classical Lagrangian for the Dirac theory is
\begin{equation}
  \mathcal{L} = \bar{\psi}(i\gamma^\mu\partial_\mu - m)\psi .
  \label{eq:3.84}
\end{equation}
The conjugate momentum to $\psi_\alpha$ is
\begin{equation}
  \pi_\alpha = \frac{\partial\mathcal{L}}{\partial\dot{\psi}_\alpha}
    = i(\psi^\dagger\gamma^0\gamma^0)_\alpha = i\psi^\dagger_\alpha .
  \label{eq:3.85}
\end{equation}
The Hamiltonian density is
\begin{equation}
  \mathcal{H} = \pi\dot{\psi} - \mathcal{L}
    = \bar{\psi}(-i\gamma^i\partial_i + m)\psi .
  \label{eq:3.86}
\end{equation}

We attempt to quantize the theory by imposing canonical \emph{commutation}
relations,
\begin{equation}
  [\psi_\alpha(\mathbf{x}), \psi_\beta^\dagger(\mathbf{y})]
    = \delta^{(3)}(\mathbf{x}-\mathbf{y})\delta_{\alpha\beta};
  \qquad
  [\psi_\alpha(\mathbf{x}), \psi_\beta(\mathbf{y})]
  = [\psi_\alpha^\dagger(\mathbf{x}), \psi_\beta^\dagger(\mathbf{y})] = 0 .
  \label{eq:3.87}
\end{equation}
We expand $\psi(x)$ and $\bar{\psi}(x)$ as Fourier integrals in terms of
creation and annihilation operators $a$, $a^\dagger$, $b$, $b^\dagger$:
\begin{equation}
  \psi(x) = \int \frac{d^3p}{(2\pi)^3}\ \frac{1}{\sqrt{2E_{\mathbf{p}}}}
    \sum_s \Bigl( a_{\mathbf{p}}^s u^s(p)e^{-ip\cdot x}
      + b_{\mathbf{p}}^{s\dagger} v^s(p)e^{+ip\cdot x} \Bigr);
  \label{eq:3.88}
\end{equation}
\begin{equation}
  \bar{\psi}(x) = \int \frac{d^3p}{(2\pi)^3}\ \frac{1}{\sqrt{2E_{\mathbf{p}}}}
    \sum_s \Bigl( b_{\mathbf{p}}^s \bar{v}^s(p)e^{-ip\cdot x}
      + a_{\mathbf{p}}^{s\dagger} \bar{u}^s(p)e^{+ip\cdot x} \Bigr).
  \label{eq:3.89}
\end{equation}
Just as for the Klein--Gordon field, we find that the commutation relations
\eqref{eq:3.87} imply
\begin{equation}
  \{a_{\mathbf{p}}^r, a_{\mathbf{q}}^{s\dagger}\}
  = \{b_{\mathbf{p}}^r, b_{\mathbf{q}}^{s\dagger}\}
    = (2\pi)^3\delta^{(3)}(\mathbf{p}-\mathbf{q})\delta^{rs},
  \label{eq:3.90}
\end{equation}
with all other commutators vanishing, provided the fields satisfy the Dirac
equation.  This is a commutation relation that leads to trouble, as we now
show.

Let us consider the propagation amplitude $\langle 0|\bar{\psi}(y)\psi(x)|0\rangle$
and its cousin $\langle 0|\psi(x)\bar{\psi}(y)|0\rangle$.  The first
amplitude receives a contribution only from the term in $\psi(x)$ with $b^\dagger$
and the term in $\bar{\psi}(y)$ with $b$; the second receives a contribution
only from the terms with $a$ and $a^\dagger$.  But for spacelike separation,
causality requires that these two amplitudes cancel:
\begin{equation}
  \langle 0|\psi_a(x)\bar{\psi}_b(y)|0\rangle
  = \langle 0|\bar{\psi}_b(y)\psi_a(x)|0\rangle
  \qquad\text{for } (x-y)^2<0 .
  \label{eq:3.91}
\end{equation}
Let us compute these amplitudes, working with the assumed commutation
relations \eqref{eq:3.87}.  We find, for example,
\begin{equation}
  \langle 0|\psi_a(x)\bar{\psi}_b(y)|0\rangle =
    \int \frac{d^3p}{(2\pi)^3}\ \frac{1}{2E_{\mathbf{p}}}
    \sum_s u_a^s(p)\bar{u}_b^s(p)\, e^{-ip\cdot(x-y)},
  \label{eq:3.92}
\end{equation}
and
\begin{equation}
  \langle 0|\bar{\psi}_b(y)\psi_a(x)|0\rangle =
    \int \frac{d^3p}{(2\pi)^3}\ \frac{1}{2E_{\mathbf{p}}}
    \sum_s v_a^s(p)\bar{v}_b^s(p)\, e^{+ip\cdot(x-y)},
  \label{eq:3.93}
\end{equation}
just as for the Klein--Gordon field.  But in the Klein--Gordon case, we got
one term of the commutator from each of these two pieces: the propagation of
a particle from $y$ to $x$ was canceled by the propagation of an antiparticle
from $x$ to $y$ outside the light-cone.  Here both terms come from the first
piece, $\langle 0|\bar{\psi}(x)\psi(y)|0\rangle$, since the second piece is
zero.  The cancellation is between positive-energy particles and
negative-energy particles, both propagating from $y$ to $x$.

This observation can actually lead us to a resolution of the negative-energy
problem.  One of the assumptions we made in quantizing the Dirac theory must
have been incorrect.  Let us therefore forget about the postulated
commutation relations \eqref{eq:3.87} and \eqref{eq:3.90}, and see whether we
can find a way for positive-energy particles to propagate in both directions.
We will also have to drop our definition of the vacuum $|0\rangle$ as the
state that is annihilated by all $a_{\mathbf{p}}^s$ and $b_{\mathbf{p}}^s$.
We will, however, retain the expressions \eqref{eq:3.88} and \eqref{eq:3.89}
for $\psi(x)$ and $\bar{\psi}(x)$ as Heisenberg operators, since if
$\psi(x)$ and $\bar{\psi}(x)$ solve the Dirac equation, they must be
decomposable into such plane-wave solutions.

First consider the propagation amplitude $\langle 0|\psi(x)\bar{\psi}(y)|0\rangle$,
which is to represent a positive-energy particle propagating from $y$ to $x$.
In this case we want the (Heisenberg) state $\bar{\psi}(y)|0\rangle$ to be
made up of only positive-energy, or negative-frequency components (since a
Heisenberg state $\psi_H = e^{iHt}\psi_S$).  Thus only the $a^\dagger$ term
of $\bar{\psi}(y)$ can contribute, which means that $a_{\mathbf{p}}^s$ must
annihilate the vacuum.  Similarly $\langle 0|\psi(x)$ can contain only
positive-frequency components.  Thus we have
\begin{equation}
  \langle 0|\psi_a(x)\bar{\psi}_b(y)|0\rangle
    = \int \frac{d^3p}{(2\pi)^3}\ \frac{1}{2E_{\mathbf{p}}}
    \sum_s u_a^s(p)\bar{u}_b^s(p)\, e^{-ip\cdot(x-y)}
    \langle 0|a_{\mathbf{p}}^s a_{\mathbf{p}}^{s\dagger}|0\rangle .
  \label{eq:3.94}
\end{equation}
We can say something about the matrix element
$\langle 0|a_{\mathbf{p}}^r a_{\mathbf{q}}^{s\dagger}|0\rangle$ even without
knowing how to interchange $a_{\mathbf{p}}^r$ and $a_{\mathbf{q}}^{s\dagger}$,
by using translational and rotational invariance.  If the ground state $|0\rangle$
is to be invariant under translations, we must have
$|0\rangle = e^{iP\cdot x}|0\rangle$.  Furthermore, since $a_{\mathbf{p}}^{r\dagger}$
creates momentum $q$, we can use Eq.~\eqref{eq:2.48} to compute
\begin{equation}
  \langle 0|a_{\mathbf{p}}^r a_{\mathbf{q}}^{s\dagger}|0\rangle
    = \langle 0|a_{\mathbf{p}}^r a_{\mathbf{q}}^{s\dagger}
      e^{iP\cdot x}|0\rangle
    = e^{i(p-q)\cdot x}\langle 0|a_{\mathbf{p}}^r a_{\mathbf{q}}^{s\dagger}|0\rangle .
\end{equation}
This says that if $\langle 0|a_{\mathbf{p}}^r a_{\mathbf{q}}^{s\dagger}|0\rangle$
is to be nonzero, $p$ must equal $q$.  Similarly, it can be shown that
rotational invariance of $|0\rangle$ implies $r = s$.  (This should be
intuitively clear, and can be checked after we discuss the angular momentum
operator later in this section.)  From these considerations we conclude that
the matrix element can be written
\begin{equation}
  \langle 0|a_{\mathbf{p}}^r a_{\mathbf{q}}^{s\dagger}|0\rangle
    = (2\pi)^3\delta^{(3)}(\mathbf{p}-\mathbf{q})\delta^{rs} \cdot A(p),
\end{equation}
where $A(p)$ is so far undetermined.  Note, however, that if the norm of a
state is always positive (as it should be in any self-respecting Hilbert
space), $A(p)$ must be greater than zero.  We can now go back to
\eqref{eq:3.94}, and write
\begin{equation}
  \langle 0|\psi_a(x)\bar{\psi}_b(y)|0\rangle
    = (i\slashed{\partial}_x + m)_{ab}\,
      \int \frac{d^3p}{(2\pi)^3}\ \frac{1}{2E_{\mathbf{p}}}\,
      e^{-ip\cdot(x-y)}\, A(p) .
  \label{eq:3.95}
\end{equation}
This expression is properly invariant under boosts only if $A(p)$ is a
Lorentz scalar, i.e., $A(p) = A(p^2)$.  Since $p^2 = m^2$, $A$ must be a
constant.  So finally we obtain
\begin{equation}
  \langle 0|\psi_a(x)\bar{\psi}_b(y)|0\rangle
    = (i\slashed{\partial}_x + m)_{ab}\,
      \int \frac{d^3p}{(2\pi)^3}\ \frac{1}{2E_{\mathbf{p}}}\,
      e^{-ip\cdot(x-y)} \cdot A .
  \label{eq:3.96}
\end{equation}
Similarly, in the amplitude $\langle 0|\bar{\psi}(y)\psi(x)|0\rangle$, we
want the only contributions to be from the positive-frequency terms of
$\bar{\psi}(y)$ and the negative-frequency terms of $\psi(x)$.  So
$a^\dagger$ still annihilates the vacuum, but $b^\dagger$ does not.  Then by
arguments identical to those given above, we have
\begin{equation}
  \langle 0|\bar{\psi}_b(y)\psi_a(x)|0\rangle
    = -(i\slashed{\partial}_x + m)_{ab}\,
      \int \frac{d^3p}{(2\pi)^3}\ \frac{1}{2E_{\mathbf{p}}}\,
      e^{+ip\cdot(x-y)} \cdot B ,
  \label{eq:3.97}
\end{equation}
where $B$ is another positive constant.  The minus sign is important; it
comes from the completeness relation \eqref{eq:3.70} for $v^s$ and the sign
of $x$ in the exponential factor.  It implies that we cannot have
$\langle 0|[\psi(x),\bar{\psi}(y)]|0\rangle = 0$ outside the light-cone: The
two terms \eqref{eq:3.96} and \eqref{eq:3.97} would indeed cancel if
$A = -B$, but this is impossible since $A$ and $B$ must both be positive.

The solution, however, is now at hand.  By setting $A = B = 1$, it is easy to
obtain (outside the light-cone)
\begin{equation}
  \langle 0|\psi_a(x)\bar{\psi}_b(y)|0\rangle
    = -\langle 0|\bar{\psi}_b(y)\psi_a(x)|0\rangle .
\end{equation}
That is, the spinor fields \emph{anticommute} at spacelike separation.  This
is enough to preserve causality, since all reasonable observables (such as
energy, charge, and particle number) are built out of an even number of
spinor fields; for any such observables $\mathcal{O}_1$ and $\mathcal{O}_2$,
we still have $[\mathcal{O}_1(x), \mathcal{O}_2(y)] = 0$ for
$(x-y)^2<0$.  And remarkably, postulating anticommutation relations for the
Dirac field solves the negative energy problem.  The equal-time
anticommutation relations will be
\begin{equation}
  \{\psi_a(\mathbf{x}), \psi_b^\dagger(\mathbf{y})\}
    = \delta^{(3)}(\mathbf{x}-\mathbf{y})\delta_{ab};
  \qquad
  \{\psi_a(\mathbf{x}), \psi_b(\mathbf{y})\}
  = \{\psi_a^\dagger(\mathbf{x}), \psi_b^\dagger(\mathbf{y})\} = 0 .
  \label{eq:3.98}
\end{equation}
We can expand $\psi(x)$ in terms of $a_{\mathbf{p}}^s$ and $b_{\mathbf{p}}^{s\dagger}$
as before (Eq.~\eqref{eq:3.88}).  The creation and annihilation operators
must now obey
\begin{equation}
  \{a_{\mathbf{p}}^r, a_{\mathbf{q}}^{s\dagger}\}
  = \{b_{\mathbf{p}}^r, b_{\mathbf{q}}^{s\dagger}\}
    = (2\pi)^3\delta^{(3)}(\mathbf{p}-\mathbf{q})\delta^{rs}
  \label{eq:3.99}
\end{equation}
(with all other anticommutators equal to zero) in order that \eqref{eq:3.98}
be satisfied.

Another computation gives the Hamiltonian,
\begin{equation}
  H = \int \frac{d^3p}{(2\pi)^3}\ E_{\mathbf{p}}
    \sum_s \bigl( a_{\mathbf{p}}^{s\dagger}a_{\mathbf{p}}^s
      - b_{\mathbf{p}}^s b_{\mathbf{p}}^{s\dagger} \bigr),
\end{equation}
which is the same as before; $b_{\mathbf{p}}^{s\dagger}$ still creates
negative energy.  However, the relation
$\{b_{\mathbf{p}}^r, b_{\mathbf{q}}^{s\dagger}\}
= (2\pi)^3\delta^{(3)}(\mathbf{p}-\mathbf{q})\delta^{rs}$ is symmetric
between $b_{\mathbf{p}}^r$ and $b_{\mathbf{q}}^{s\dagger}$.  So let us simply
redefine
\begin{equation}
  b_{\mathbf{p}}^s = B_{\mathbf{p}}^s,\qquad
  b_{\mathbf{p}}^{s\dagger} = B_{\mathbf{p}}^{s\dagger} .
  \label{eq:3.100}
\end{equation}
These of course obey exactly the same anticommutation relations, but now the
second term in the Hamiltonian is
\begin{equation}
  -E_{\mathbf{p}}b_{\mathbf{p}}^s b_{\mathbf{p}}^{s\dagger}
    = +E_{\mathbf{p}}b_{\mathbf{p}}^{s\dagger}b_{\mathbf{p}}^s
      - (\text{const}).
\end{equation}
If we choose $|0\rangle$ to be the state that is annihilated by
$a_{\mathbf{p}}^s$ and $b_{\mathbf{p}}^s$, then all excitations of $|0\rangle$
have positive energy.

What happened?  To better understand this trick, let us abandon the field
theory for a moment and consider a theory with a single pair of $b$ and
$b^\dagger$ operators obeying $\{b,b^\dagger\} = 1$ and $\{b,b\} =
\{b^\dagger,b^\dagger\} = 0$.  Choose a state $|0\rangle$ such that $b|0\rangle = 0$.
Then $b^\dagger|0\rangle$ is a new state; call it $|1\rangle$.  This state
satisfies $b|1\rangle = |0\rangle$ and $b^\dagger|1\rangle = 0$.  So $b$ and
$b^\dagger$ act on a Hilbert space of only two states, $|0\rangle$ and
$|1\rangle$.  We might say that $|0\rangle$ represents an ``empty'' state, and
that $b^\dagger$ ``fills'' the state.  But we could equally well call
$|1\rangle$ the empty state and say that $b$ --- fills it.  The two
descriptions are completely equivalent, until we specify some observable that
allows us to distinguish the states physically.  In our case the correct
choice is to take the state of lower energy to be the empty one.  And it is
less confusing to put the dagger on the operator that creates positive
energy.  That is exactly what we have done.

Note, by the way, that since $(b^\dagger)^2 = 0$, the state cannot be filled
twice.  More generally, the anticommutation relations imply that any
multiparticle state is antisymmetric under the interchange of two particles:
$a_{\mathbf{p}}^{s\dagger}a_{\mathbf{q}}^{r\dagger}|0\rangle =
-a_{\mathbf{q}}^{r\dagger}a_{\mathbf{p}}^{s\dagger}|0\rangle$.  Thus we
conclude that if the ladder operators obey anticommutation relations, the
corresponding particles obey Fermi--Dirac statistics.

We have just shown that in order to insure that the vacuum has only
positive-energy excitations, we must quantize the Dirac field with
anticommutation relations; under these conditions the particles associated
with the Dirac field obey Fermi--Dirac statistics.  This conclusion is part
of a more general result, first derived by Pauli.%
\footnote{W.~Pauli, \emph{Phys.\ Rev.} \textbf{58}, 716 (1940), reprinted in
Schwinger (1958).  A rigorous treatment is given by R.~F.~Streater and
A.~S.~Wightman, \emph{PCT, Spin and Statistics, and All That}
(Benjamin/Cummings, Reading, Mass., 1964).}: Lorentz invariance, positive
energies, positive norms, and causality together imply that particles of
integer spin obey Bose--Einstein statistics, while particles of half-odd-integer
spin obey Fermi--Dirac statistics.

\subsection{The Quantized Dirac Field}

Let us now summarize the results of the quantized Dirac theory in a
systematic way.  Since the dust has settled, we should clean up our notation:
From now on we will write $b_{\mathbf{p}}^s$ (the operator that lowers the
energy of a state) simply as $b_{\mathbf{p}}^s$, and $b_{\mathbf{p}}^{s\dagger}$
as $b_{\mathbf{p}}^{s\dagger}$.  All the expressions we will need in our later
work are listed below; corresponding expressions above, where they differ,
should be forgotten.

First we write the field operators:
\begin{equation}
  \psi(x) = \int \frac{d^3p}{(2\pi)^3}\ \frac{1}{\sqrt{2E_{\mathbf{p}}}}
    \sum_s \Bigl( a_{\mathbf{p}}^s u^s(p)e^{-ip\cdot x}
      + b_{\mathbf{p}}^{s\dagger} v^s(p)e^{+ip\cdot x} \Bigr);
  \label{eq:3.101}
\end{equation}
\begin{equation}
  \bar{\psi}(x) = \int \frac{d^3p}{(2\pi)^3}\ \frac{1}{\sqrt{2E_{\mathbf{p}}}}
    \sum_s \Bigl( b_{\mathbf{p}}^s \bar{v}^s(p)e^{-ip\cdot x}
      + a_{\mathbf{p}}^{s\dagger} \bar{u}^s(p)e^{+ip\cdot x} \Bigr).
  \label{eq:3.102}
\end{equation}
The creation and annihilation operators obey the anticommutation rules
\begin{equation}
  \{a_{\mathbf{p}}^r, a_{\mathbf{q}}^{s\dagger}\}
  = \{b_{\mathbf{p}}^r, b_{\mathbf{q}}^{s\dagger}\}
    = (2\pi)^3\delta^{(3)}(\mathbf{p}-\mathbf{q})\delta^{rs},
  \label{eq:3.103}
\end{equation}
with all other anticommutators equal to zero.  The equal-time anticommutation
relations for $\psi$ and $\psi^\dagger$ are then
\begin{equation}
  \{\psi_a(\mathbf{x}), \psi_b^\dagger(\mathbf{y})\}
    = \delta^{(3)}(\mathbf{x}-\mathbf{y})\delta_{ab};
  \qquad
  \{\psi_a(\mathbf{x}), \psi_b(\mathbf{y})\}
  = \{\psi_a^\dagger(\mathbf{x}), \psi_b^\dagger(\mathbf{y})\} = 0 .
  \label{eq:3.104}
\end{equation}
The vacuum $|0\rangle$ is defined to be the state such that
\begin{equation}
  a_{\mathbf{p}}^s|0\rangle = b_{\mathbf{p}}^s|0\rangle = 0 .
  \label{eq:3.105}
\end{equation}
The Hamiltonian can be written
\begin{equation}
  H = \int \frac{d^3p}{(2\pi)^3}\ \sum_s E_{\mathbf{p}}
    \Bigl( a_{\mathbf{p}}^{s\dagger}a_{\mathbf{p}}^s
      + b_{\mathbf{p}}^{s\dagger}b_{\mathbf{p}}^s \Bigr),
  \label{eq:3.106}
\end{equation}
where we have dropped the infinite constant term that comes from
anticommuting $b_{\mathbf{p}}^s$ and $b_{\mathbf{p}}^{s\dagger}$.  From this we
see that the vacuum is the state of lowest energy, as desired.  The momentum
operator is
\begin{equation}
  \mathbf{P} = \int d^3x\ \psi^\dagger(-i\nabla)\psi
    = \int \frac{d^3p}{(2\pi)^3}\ \sum_s \mathbf{p}
      \bigl( a_{\mathbf{p}}^{s\dagger}a_{\mathbf{p}}^s
        + b_{\mathbf{p}}^{s\dagger}b_{\mathbf{p}}^s \bigr).
  \label{eq:3.107}
\end{equation}
Thus both $a_{\mathbf{p}}^{s\dagger}$ and $b_{\mathbf{p}}^{s\dagger}$ create
particles with energy $+E_{\mathbf{p}}$ and momentum $\mathbf{p}$.  We will
refer to the particles created by $a_{\mathbf{p}}^{s\dagger}$ as
\emph{fermions} and to those created by $b_{\mathbf{p}}^{s\dagger}$ as
\emph{antifermions}.

The one-particle states
\begin{equation}
  |\mathbf{p}, s\rangle = \sqrt{2E_{\mathbf{p}}}\,a_{\mathbf{p}}^{s\dagger}|0\rangle
  \label{eq:3.108}
\end{equation}
are defined so that their inner product
\begin{equation}
  \langle\mathbf{p}, r|\mathbf{q}, s\rangle
    = 2E_{\mathbf{p}}\,(2\pi)^3\delta^{(3)}(\mathbf{p}-\mathbf{q})\delta^{rs}
  \label{eq:3.109}
\end{equation}
is Lorentz invariant.  This implies that the operator $U(\Lambda)$ that
implements Lorentz transformations on the states of the Hilbert space is
unitary, even though for boosts, $\Lambda_{1/2}$ is not unitary.

It will be reassuring to do a consistency check, to see that $U(\Lambda)$
implements the right transformation on $\psi(x)$.  So calculate
\begin{equation}
  U\psi(x)U^{-1} = U \int \frac{d^3p}{(2\pi)^3}\ \frac{1}{\sqrt{2E_{\mathbf{p}}}}
    \sum_s \Bigl( a_{\mathbf{p}}^s u^s(p)e^{-ip\cdot x}
      + b_{\mathbf{p}}^{s\dagger} v^s(p)e^{ip\cdot x} \Bigr) U^{-1}.
  \label{eq:3.110}
\end{equation}
We can concentrate on the first term; the second is completely analogous.
Equation \eqref{eq:3.108} implies that $a_{\mathbf{p}}^{s\dagger}$ transforms
according to
\begin{equation}
  U(\Lambda)a_{\mathbf{p}}^{s\dagger}U^{-1}(\Lambda)
    = a_{\Lambda\mathbf{p}}^{s\dagger},
  \label{eq:3.111}
\end{equation}
assuming that the axis of spin quantization is parallel to the boost or
rotation axis.  To use this relation to evaluate \eqref{eq:3.110}, rewrite
the integral as
\begin{equation}
  \int \frac{d^3p}{(2\pi)^3}\ \frac{1}{\sqrt{2E_{\mathbf{p}}}}\ a_{\mathbf{p}}^s
    = \int \frac{d^3p}{(2\pi)^3}\ \frac{1}{\sqrt{2E_{\mathbf{p}}}}\,
      u^s(p)\, e^{-ip\cdot x}
    = \int \frac{d^3p}{(2\pi)^3}\ \frac{1}{2E_{\mathbf{p}}}\,
      u^s(p)\, e^{-ip\cdot x}\, \sqrt{2E_{\mathbf{p}}} .
\end{equation}
The second factor is transformed in a simple way by $U$, and the first is a
Lorentz-invariant integral.  Thus, if we apply \eqref{eq:3.111} and make the
substitution $p = \Lambda p$, Eq.~\eqref{eq:3.110} becomes
\begin{equation}
  U(\Lambda)\psi(x)U^{-1}(\Lambda) =
    \int \frac{d^3p}{(2\pi)^3}\ \frac{1}{\sqrt{2E_{\mathbf{p}}}}\,
    \sum_s a_{\mathbf{p}}^{s\dagger}
    u^s(\Lambda^{-1}p)\, e^{-ip\cdot \Lambda x} + \cdots .
\end{equation}
But $u^s(\Lambda^{-1}p) = \Lambda_{1/2}u^s(p)$, so indeed we have
\begin{equation}
  U(\Lambda)\psi(x)U^{-1}(\Lambda) = \Lambda_{1/2}\psi(\Lambda x),
  \label{eq:3.112}
\end{equation}
as advertised.

There is one more important conserved quantity in the Dirac theory.  In
Section~3.4 we saw that the current $j^\mu = \bar{\psi}\gamma^\mu\psi$ is
conserved.  The charge associated with this current is
\begin{equation}
  Q = \int d^3x\ \psi^\dagger(x)\psi(x)
    = \int \frac{d^3p}{(2\pi)^3}\ \sum_s
      \bigl( a_{\mathbf{p}}^{s\dagger}a_{\mathbf{p}}^s
        - b_{\mathbf{p}}^{s\dagger}b_{\mathbf{p}}^s \bigr),
  \label{eq:3.113}
\end{equation}
or, if we ignore another infinite constant,
\begin{equation}
  Q = \int \frac{d^3p}{(2\pi)^3}\ \sum_s
    \bigl( a_{\mathbf{p}}^{s\dagger}a_{\mathbf{p}}^s
      - b_{\mathbf{p}}^s b_{\mathbf{p}}^{s\dagger} \bigr)
    = \int \frac{d^3p}{(2\pi)^3}\ \sum_s
      \bigl( a_{\mathbf{p}}^{s\dagger}a_{\mathbf{p}}^s
        - b_{\mathbf{p}}^{s\dagger}b_{\mathbf{p}}^s \bigr).
  \label{eq:3.114}
\end{equation}
So $a_{\mathbf{p}}^{s\dagger}$ creates fermions with charge $+1$, while
$b_{\mathbf{p}}^{s\dagger}$ creates antifermions with charge $-1$.  When we
couple the Dirac field to the electromagnetic field, we will see that $Q$ is
none other than the electric charge (up to a constant factor that depends on
which type of particle we wish to describe; e.g., for electrons, the electric
charge is $Qe$).

In Quantum Electrodynamics we will use the spinor field $\psi$ to describe
electrons and positrons.  The particles created by $a_{\mathbf{p}}^{s\dagger}$
are electrons; they have energy $E_{\mathbf{p}}$, momentum $\mathbf{p}$, spin
$1/2$ with polarization appropriate to $u^s$, and charge $+1$ (in units of
$e$).  The particles created by $b_{\mathbf{p}}^{s\dagger}$ are positrons;
they have energy $E_{\mathbf{p}}$, momentum $\mathbf{p}$, spin $1/2$ with
polarization opposite to that of $u^s$, and charge $-1$.  The state
$\psi_a(x)|0\rangle$ contains a positron at position $x$, whose polarization
corresponds to the spinor component chosen.  Similarly, $\bar{\psi}_a(x)|0\rangle$
is a state of one electron at position $x$.

\subsection{The Dirac Propagator}

Calculating propagation amplitudes for the Dirac field is by now a
straightforward exercise:
\begin{align}
  \langle 0|\psi_a(x)\bar{\psi}_b(y)|0\rangle
  &= \int \frac{d^3p}{(2\pi)^3}\ \frac{1}{2E_{\mathbf{p}}}
    (i\slashed{\partial}_x + m)_{ab}\, e^{-ip\cdot(x-y)}
  = (i\slashed{\partial}_x + m)_{ab}\,
    \int \frac{d^3p}{(2\pi)^3}\ \frac{1}{2E_{\mathbf{p}}}\,
    e^{-ip\cdot(x-y)} ,
  \label{eq:3.115}\\
  \langle 0|\bar{\psi}_b(y)\psi_a(x)|0\rangle
  &= (i\slashed{\partial}_x + m)_{ab}\,
    \int \frac{d^3p}{(2\pi)^3}\ \frac{1}{2E_{\mathbf{p}}}\,
    e^{+ip\cdot(x-y)} .
  \label{eq:3.116}
\end{align}
Just as we did for the Klein--Gordon equation, we can construct Green's
functions for the Dirac equation obeying various boundary conditions.  For
example, the retarded Green's function is
\begin{equation}
  S_R(x-y) = \theta(x^0-y^0)\ \langle 0|\{\psi_a(x), \bar{\psi}_b(y)\}|0\rangle .
  \label{eq:3.117}
\end{equation}
It is easy to verify that
\begin{equation}
  S_R(x-y) = (i\slashed{\partial}_x + m)\,D_R(x-y),
  \label{eq:3.118}
\end{equation}
since on the right-hand side the term involving
$\partial_0\theta(x^0-y^0)$ vanishes.  Using \eqref{eq:3.118} and the fact
that $\slashed{\partial}\slashed{\partial} = \partial^2$, we see that $S_R$
is a Green's function of the Dirac operator:
\begin{equation}
  (i\slashed{\partial}_x - m)\,S_R(x-y) = i\delta^{(4)}(x-y)\cdot
    \mathbb{1}_{4\times4} .
  \label{eq:3.119}
\end{equation}
The Green's function of the Dirac operator can also be found by Fourier
transformation.  Expanding $S_R(x-y)$ as a Fourier integral and acting on
both sides with $(i\slashed{\partial}_x - m)$, we find
\begin{equation}
  iS_R(x-y) = \int \frac{d^4p}{(2\pi)^4}\ (\slashed{p} - m)\,
    e^{-ip\cdot(x-y)}\, S_R(p),
  \label{eq:3.120}
\end{equation}
and hence
\begin{equation}
  S_R(p) = \frac{i}{\slashed{p} - m}
    = \frac{i(\slashed{p} + m)}{p^2 - m^2} .
  \label{eq:3.121}
\end{equation}
To obtain the retarded Green's function, we must evaluate the $p^0$ integral
in \eqref{eq:3.121} along the contour shown on page~30.  For $x^0>y^0$ we
close the contour below, picking up both poles to obtain the sum of
\eqref{eq:3.115} and \eqref{eq:3.116}.  For $x^0<y^0$ we close the contour
above and get zero.

The Green's function with Feynman boundary conditions is defined by the
contour shown on page~31:
\begin{equation}
  S_F(x-y) = \int \frac{d^4p}{(2\pi)^4}\
    \frac{i(\slashed{p} + m)}{p^2 - m^2 + i\epsilon}\ e^{-ip\cdot(x-y)}
  = \begin{cases}
      \langle 0|\psi(x)\bar{\psi}(y)|0\rangle & \text{for } x^0>y^0
        \text{ (close contour below)}\\
      -\langle 0|\bar{\psi}(y)\psi(x)|0\rangle & \text{for } x^0<y^0
        \text{ (close contour above)}
    \end{cases}
  = \langle 0|T\psi(x)\bar{\psi}(y)|0\rangle ,
  \label{eq:3.122}
\end{equation}
where we have chosen to define the time-ordered product of spinor fields with
an additional minus sign when the operators are interchanged.  This minus
sign is extremely important in the quantum field theory of fermions; we will
meet it again in Section~4.7.

As with the Klein--Gordon theory, the expression \eqref{eq:3.122} for the
Feynman propagator is the most useful result of this chapter.  When we do
perturbative calculations with Feynman diagrams, we will associate the factor
$S_F(p)$ with each internal fermion line.

\section{Discrete Symmetries of the Dirac Theory}

In the last section we discussed the implementation of continuous Lorentz
transformations on the Hilbert space of the Dirac theory.  We found that for
each transformation $\Lambda$ there was a unitary operator $U(\Lambda)$, which
induced the correct transformation on the fields:
\begin{equation}
  U(\Lambda)\psi(x)U^{-1}(\Lambda) = \Lambda_{1/2}\psi(\Lambda x).
  \label{eq:3.123}
\end{equation}
In this section we will discuss the analogous operators that implement various
discrete symmetries on the Dirac field.

In addition to continuous Lorentz transformations, there are two other
spacetime operations that are potential symmetries of the Lagrangian: parity
and time reversal.  Parity, denoted by $P$, sends $(t,\mathbf{x})\to(t,-\mathbf{x})$,
reversing the handedness of space.  Time reversal, denoted by $T$, sends
$(t,\mathbf{x})\to(-t,\mathbf{x})$, interchanging the forward and backward
light-cones.  Neither of these operations can be achieved by a continuous
Lorentz transformation starting from the identity.  Both, however, preserve
the Minkowski interval $x^2 = t^2 - \mathbf{x}^2$.  In standard terminology,
the continuous Lorentz transformations are referred to as the proper,
orthochronous Lorentz group, $L_+$.  Then the full Lorentz group breaks up
into four disconnected subsets, as shown below:
\begin{equation}
  \begin{array}{ll}
    L_+^\uparrow = L_+ & \text{``proper''}\\
    L_- = PL_+ & \text{``improper''}\\
    L_+^\downarrow = TL_+ & \text{``nonorthochronous''}\\
    L_-^\uparrow = PL_+^\downarrow = TL_+^\uparrow &
      \text{``improper nonorthochronous''}
  \end{array}
\end{equation}
At the same time that we discuss $P$ and $T$, it will be convenient to
discuss a third (non-spacetime) discrete operation: charge conjugation,
denoted by $C$.  Under this operation, particles and antiparticles are
interchanged.

Although any relativistic field theory must be invariant under $L_+^\uparrow$,
it need not be invariant under $P$, $T$, or $C$.  What is the status of these
symmetry operations in the real world?  From experiment, we know that three
of the forces of Nature---the gravitational, electromagnetic, and strong
interactions---are symmetric with respect to $P$, $C$, and $T$.  The weak
interactions violate $C$ and $P$ separately, but preserve $CP$ and $T$.  But
certain rare processes (all so far observed involve neutral $K$ mesons) also
show $CP$ and $T$ violation.  All observations indicate that the combination
$CPT$ is a perfect symmetry of Nature.

The currently accepted theoretical model of the weak interactions is the
Glashow--Weinberg--Salam gauge theory, described in Chapter~20.  This theory
violates $C$ and $P$ in the strongest possible way.  It is actually a surprise
(though not quite an accident) that $C$ and $P$ happen to be quite good
symmetries in the most readily observable processes.  On the other hand, no
one knows a really beautiful theory that violates $CP$.  In the current
theory, when there are three (or more) fermion generations, there is room for
a parameter that, if nonzero, causes $CP$ violation.  But the value of this
parameter is no better understood than the value of the electron mass; the
physical origin of $CP$ violation remains a mystery.  We will discuss this
question further in Section~20.3.

\subsection{Parity}

With this introduction, let us now discuss the action of $P$, $T$, and $C$ on
Dirac particles and fields.  First consider parity.  The operator $P$ should
reverse the momentum of a particle without flipping its spin:
\begin{equation}
  \mathbf{p} \to -\mathbf{p}\qquad \text{(mirror)}.
\end{equation}
Mathematically, this means that $P$ should be implemented by a unitary
operator (properly called $U(P)$, but we'll just call it $P$) which, for
example, transforms the state $a_{\mathbf{p}}^{s\dagger}|0\rangle$ into
$a_{-\mathbf{p}}^{s\dagger}|0\rangle$.  In other words, we want
\begin{equation}
  P a_{\mathbf{p}}^s P = \eta_a a_{-\mathbf{p}}^s
  \qquad\text{and}\qquad
  P b_{\mathbf{p}}^s P = \eta_b b_{-\mathbf{p}}^s,
  \label{eq:3.124}
\end{equation}
where $\eta_a$ and $\eta_b$ are possible phases.  These phases are restricted
by the condition that two applications of the parity operator should return
observables to their original values.  Since observables are built from an
even number of fermion operators, this requires $\eta_a^2, \eta_b^2 = \pm 1$.

Just as a continuous Lorentz transformation is implemented on the Dirac field
as the $4\times4$ constant matrix $\Lambda_{1/2}$, the parity transformation
should also be represented by a $4\times4$ constant matrix.  To find this
matrix, and to determine $\eta_a$ and $\eta_b$, we compute the action of $P$
on $\psi(x)$.  Using \eqref{eq:3.124}, we have
\begin{equation}
  P\psi(x)P = \int \frac{d^3p}{(2\pi)^3}\ \frac{1}{\sqrt{2E_{\mathbf{p}}}}
    \sum_s \Bigl( \eta_a a_{-\mathbf{p}}^s u^s(p)e^{-ip\cdot x}
      + \eta_b^{*} b_{\mathbf{p}}^{s\dagger} v^s(p)e^{ip\cdot x} \Bigr).
  \label{eq:3.125}
\end{equation}
Now change variables to $\mathbf{p}' = (p^0, -\mathbf{p})$.  Note that
$p\cdot x = p'\cdot(t,-\mathbf{x})$.  Also
\begin{equation}
  u^s(p) = \gamma^0 u^s(p'),\qquad
  v^s(p) = -\gamma^0 v^s(p').
\end{equation}
So \eqref{eq:3.125} becomes
\begin{equation}
  P\psi(x)P = \int \frac{d^3p}{(2\pi)^3}\ \frac{1}{\sqrt{2E_{\mathbf{p}}}}
    \sum_s \Bigl( \eta_a a_{\mathbf{p}}^s \gamma^0 u^s(p)e^{-ip\cdot
      (t,-\mathbf{x})} - \eta_b^{*} b_{\mathbf{p}}^{s\dagger} \gamma^0 v^s(p)
      e^{ip\cdot(t,-\mathbf{x})} \Bigr).
\end{equation}
This should equal some constant matrix times $\psi(t,-\mathbf{x})$, and
indeed it works if we make $\eta_b^{*} = -\eta_a$.  This implies
\begin{equation}
  \eta_a\eta_b = -\eta_a^2 = -1 .
  \label{eq:3.126}
\end{equation}
Thus we have the parity transformation of $\psi(x)$ in its final form,
\begin{equation}
  P\psi(t,\mathbf{x})P = \eta_a\gamma^0\psi(t,-\mathbf{x}) .
  \label{eq:3.127}
\end{equation}
It will be very important (for example, in writing down Lagrangians) to know
how the various Dirac field bilinears transform under parity.  Recall that
the five bilinears are
\begin{equation}
  \bar{\psi}\psi,\qquad \bar{\psi}\gamma^\mu\psi,\qquad
  \bar{\psi}\sigma^{\mu\nu}\psi,\qquad
  \bar{\psi}\gamma^\mu\gamma^5\psi,\qquad
  \bar{\psi}i\gamma^5\psi .
\end{equation}
Under parity, one finds
\begin{align}
  P\bar{\psi}\psi\,P &= +\bar{\psi}\psi(t,-\mathbf{x});\\
  P\bar{\psi}\gamma^\mu\psi\,P &=
    +(-1)^\mu\bar{\psi}\gamma^\mu\psi(t,-\mathbf{x});\\
  P\bar{\psi}\sigma^{\mu\nu}\psi\,P &=
    (-1)^\mu(-1)^\nu\bar{\psi}\sigma^{\mu\nu}\psi(t,-\mathbf{x});\\
  P\bar{\psi}\gamma^\mu\gamma^5\psi\,P &=
    -(-1)^\mu\bar{\psi}\gamma^\mu\gamma^5\psi(t,-\mathbf{x});\\
  P\bar{\psi}i\gamma^5\psi\,P &=
    -\bar{\psi}i\gamma^5\psi(t,-\mathbf{x}),
\end{align}
where $(-1)^\mu = 1$ for $\mu = 0$ and $(-1)^\mu = -1$ for
$\mu = 1, 2, 3$.

\subsection{Time Reversal}

We now turn to the more subtle case of time reversal.  We will not describe
the full treatment given in the original literature; rather, we will confine
ourselves to a short presentation of the results, which will be useful in
understanding the action of $T$ on the Dirac field.

The time-reversal operation $T$ should change the direction of time and, in
a sense, ``invert'' the quantum-mechanical evolution.  A complete description
of the action of $T$ on a quantum mechanical system requires that we treat
$T$ as an \emph{antilinear} (or antiunitary) operator.  Since this is
probably a new concept, let us explain it briefly.  A unitary operator $U$
acts on a linear superposition of states as
\begin{equation}
  U(a|A\rangle + b|B\rangle) = aU|A\rangle + bU|B\rangle .
\end{equation}
An antilinear operator $\Theta$, on the other hand, acts as
\begin{equation}
  \Theta(a|A\rangle + b|B\rangle) = a^{*}\Theta|A\rangle + b^{*}\Theta|B\rangle .
\end{equation}
Time reversal is antilinear because one expects, for example, that
$\langle B|A\rangle$ (a complex number) is transformed to its complex
conjugate under $T$: $T\langle B|A\rangle = \langle B|A\rangle^{*}$.
Moreover, if $H$ is the Hamiltonian of a time-reversal-invariant system, so
that $T$ commutes with $H$, then the time-evolution operator $e^{-iHt}$
satisfies $Te^{-iHt} = e^{iHt}T$, which is the statement that $T$ reverses
the direction of time.

It is convenient to parametrize the antilinear transformation $T$ of a
Dirac field as follows.  Since $T$ reverses momenta and spins, we expect
\begin{equation}
  T\psi(t,\mathbf{x})T = \eta_T\, C\,\psi^{*}(-t,\mathbf{x})
\end{equation}
for some phase $\eta_T$ and some $4\times4$ matrix $C$ (not to be confused
with the charge conjugation operator; here $C$ is called the time-reversal
matrix).  Actually, to keep the notation simple, let us define
\begin{equation}
  T\psi(t,\mathbf{x})T = \gamma^1\gamma^3\psi^{*}(-t,\mathbf{x})
  \label{eq:3.128}
\end{equation}
(an overall phase has been suppressed).  Using the notation of Eq.~\eqref{eq:3.67},
we define the time reversal transformation of fermion annihilation operators
as follows:
\begin{equation}
  T a_{\mathbf{p}}^s T = a_{-\mathbf{p}}^s,\qquad
  T b_{\mathbf{p}}^s T = b_{-\mathbf{p}}^s .
  \label{eq:3.129}
\end{equation}
(An additional overall phase would have no effect on the rest of our
discussion and is omitted for simplicity.)  Relations \eqref{eq:3.129} allow
us to compute the action of $T$ on the fermion field $\psi(x)$:
\begin{align}
  T\psi(t,\mathbf{x})T &= \int \frac{d^3p}{(2\pi)^3}\ \frac{1}{\sqrt{2E_{\mathbf{p}}}}
    \sum_s \Bigl( a_{-\mathbf{p}}^s u^{s*}(p)e^{+ip\cdot x}
      + b_{-\mathbf{p}}^{s\dagger} v^{s*}(p)e^{-ip\cdot x} \Bigr)\\
  &= \gamma^1\gamma^3\,\psi(-t,\mathbf{x}) .
  \label{eq:3.130}
\end{align}
In the last step we used $p\cdot(t,-\mathbf{x}) = -p\cdot(-t,\mathbf{x})$.
Just as for parity, we have derived a simple transformation law for the
fermion field $\psi(x)$.  The relative minus sign in the transformation laws
for particle and antiparticle is present here as well, implicit in the
twice-flipped spinor in $v^{-s}$.

Now we can check the action of $T$ on the various bilinears.  First we need
\begin{equation}
  T\bar{\psi}T = (T\psi T)^{\dagger}\gamma^0
    = \psi^{\dagger}(-t,\mathbf{x})[\gamma^1\gamma^3]^{\dagger}\gamma^0
    = \bar{\psi}(-t,\mathbf{x})[-\gamma^1\gamma^3] .
  \label{eq:3.131}
\end{equation}
Then the transformation of the scalar bilinear is
\begin{equation}
  T\bar{\psi}\psi\,T = \bar{\psi}(-\gamma^1\gamma^3)(\gamma^1\gamma^3)\psi(-t,\mathbf{x})
    = \bar{\psi}\psi(-t,\mathbf{x}) .
  \label{eq:3.132}
\end{equation}
The pseudo-scalar acquires an extra minus sign when $T$ goes through the $i$:
\begin{equation}
  Ti\bar{\psi}\gamma^5\psi\,T = -i\bar{\psi}(-\gamma^1\gamma^3)\gamma^5
    (\gamma^1\gamma^3)\psi = -i\bar{\psi}\gamma^5\psi(-t,\mathbf{x}).
\end{equation}
For the vector, we must separately compute each of the four cases
$\mu = 0, 1, 2, 3$.  After a bit of work you should find
\begin{equation}
  T\bar{\psi}\gamma^\mu\psi\,T =
    \bar{\psi}(-\gamma^1\gamma^3)(\gamma^\mu)^{*}(\gamma^1\gamma^3)\psi
    = \begin{cases}
        +\bar{\psi}\gamma^\mu\psi(-t,\mathbf{x}) & \text{for } \mu = 0;\\
        -\bar{\psi}\gamma^\mu\psi(-t,\mathbf{x}) & \text{for } \mu = 1, 2, 3 .
      \end{cases}
  \label{eq:3.133}
\end{equation}
This is exactly the transformation property we want for vectors such as the
current density.  You can verify that the pseudo-vector transforms in exactly
the same way under time-reversal.

\subsection{Charge Conjugation}

The last of the three discrete symmetries is the particle-antiparticle
symmetry $C$.  There will be no problem in implementing $C$ as a unitary
linear operator.  Charge conjugation is conventionally defined to take a
fermion with a given spin orientation into an antifermion with the same spin
orientation.  Thus, a convenient choice for the transformation of fermion
annihilation operators is
\begin{equation}
  C a_{\mathbf{p}}^s C = b_{\mathbf{p}}^s;\qquad
  C b_{\mathbf{p}}^s C = a_{\mathbf{p}}^s .
  \label{eq:3.134}
\end{equation}
Again, we ignore possible additional phases for simplicity.

Next we want to work out the action of $C$ on $\psi(x)$.  First we need a
relation between $v^s(p)$ and $u^s(p)$.  Using \eqref{eq:3.67}, and
\eqref{eq:3.55},
\begin{equation}
  v^s(p) = \begin{pmatrix}
    \sqrt{p\cdot\sigma}\,\eta^s\\
    \sqrt{p\cdot\bar{\sigma}}\,\eta^s
  \end{pmatrix}
  = \begin{pmatrix}
    -i\sigma^2\sqrt{p\cdot\bar{\sigma}}\,\xi^{s*}\\
    i\sigma^2\sqrt{p\cdot\sigma}\,\xi^{s*}
  \end{pmatrix}
  = -i\gamma^2\begin{pmatrix}
    \sqrt{p\cdot\sigma}\,\xi^{s*}\\
    \sqrt{p\cdot\bar{\sigma}}\,\xi^{s*}
  \end{pmatrix},
\end{equation}
where we have used $v^s(p) = (u^s(p))^{c}$, the charge conjugate.  That is,
\begin{equation}
  u^s(p) = -i\gamma^2(v^s(p))^{*},\qquad
  v^s(p) = -i\gamma^2(u^s(p))^{*}.
  \label{eq:3.135}
\end{equation}
If we substitute \eqref{eq:3.135} into the expression for the fermion field
operator, and then transform this operator with $C$, we find
\begin{equation}
  C\psi(x)C = -i\gamma^2\psi^{*}(x) = -i\gamma^2(\psi^{\dagger}\gamma^0)^T
    = -i(\bar{\psi}\gamma^0\gamma^2)^T .
  \label{eq:3.136}
\end{equation}
Note that $C$ is a linear unitary operator, even though it takes $\psi \to
\psi^{c}$.

Once again, we would like to know how $C$ acts on fermion bilinears.  First
we need
\begin{equation}
  C\bar{\psi}(x)C = C\psi^{\dagger}(x)C\gamma^0
    = (-i\gamma^2\psi^{*})^{\dagger}\gamma^0
    = (-i\gamma^0\gamma^2\psi)^{\dagger} .
  \label{eq:3.137}
\end{equation}
Working out the transformations of bilinears is a bit tricky, and it helps to
write in spinor indices.  For the scalar,
\begin{align}
  C\bar{\psi}\psi\,C &= (-i\bar{\psi}\gamma^0\gamma^2)^T(-i\gamma^2\psi^{*})
    = -\psi^{\dagger}\gamma^2\gamma^0\gamma^2\psi^{*}
    = \psi^{\dagger}\gamma^0\psi
    = +\bar{\psi}\psi .
  \label{eq:3.138}
\end{align}
(The minus sign in the third step is from fermion anticommutation.)  The
pseudo-scalar is no more difficult:
\begin{equation}
  Ci\bar{\psi}\gamma^5\psi\,C = +i\bar{\psi}\gamma^5\psi .
  \label{eq:3.139}
\end{equation}
We must do each component of the vector and pseudo-vector separately.  Noting
that $\gamma^0$ and $\gamma^2$ are symmetric matrices while $\gamma^1$ and
$\gamma^3$ are antisymmetric, we eventually find
\begin{equation}
  C\bar{\psi}\gamma^\mu\psi\,C = -\bar{\psi}\gamma^\mu\psi;
  \label{eq:3.140}
\end{equation}
\begin{equation}
  C\bar{\psi}\gamma^\mu\gamma^5\psi\,C = +\bar{\psi}\gamma^\mu\gamma^5\psi .
  \label{eq:3.141}
\end{equation}
Although the operator $C$ interchanges $\psi$ and $\bar{\psi}$, it does not
actually change the order of the creation and annihilation operators.  Thus,
if $\bar{\psi}\gamma^0\psi$ is defined to subtract the infinite constant
noted above Eq.~\eqref{eq:3.113}, this constant does not reappear in the
process of conjugation by $C$.

\subsection{Summary of $C$, $P$, and $T$}

The transformation properties of the various fermion bilinears under $C$,
$P$, and $T$ are summarized in the table below.  Here we use the shorthand
$(-1)^\mu = 1$ for $\mu = 0$ and $(-1)^\mu = -1$ for
$\mu = 1, 2, 3$.
\begin{equation}
  \begin{array}{c|ccccc}
    & \bar{\psi}\psi & \bar{\psi}\gamma^5\psi
      & \bar{\psi}\gamma^\mu\psi & \bar{\psi}\gamma^\mu\gamma^5\psi
      & \bar{\psi}\sigma^{\mu\nu}\psi\\
    \hline
    P & +1 & -1 & (-1)^\mu & -(-1)^\mu & (-1)^\mu(-1)^\nu\\
    T & +1 & -1 & (-1)^\mu & (-1)^\mu & -(-1)^\mu(-1)^\nu\\
    C & +1 & +1 & -1 & +1 & -1\\
    CPT & +1 & +1 & -1 & -1 & +1
  \end{array}
  \label{eq:3.142}
\end{equation}
We have included the transformation properties of the tensor bilinear (see
Problem~3.7), and also of the derivative operator.

Notice first that the free Dirac Lagrangian
$\mathcal{L}_0 = \bar{\psi}(i\slashed{\partial} - m)\psi$ is invariant under
$C$, $P$, and $T$ separately.  We can build more general quantum systems that
violate any of these symmetries by adding to $\mathcal{L}_0$ some
perturbation $\delta\mathcal{L}$.  But $\delta\mathcal{L}$ must be a Lorentz
scalar, and the last line of the table shows that all Lorentz scalar
combinations of $\psi$ and $\bar{\psi}$ are invariant under the combined
symmetry $CPT$.  Actually, it is quite generally true that one cannot build a
Lorentz-invariant quantum field theory with a Hermitian Hamiltonian that
violates $CPT$.%
\footnote{This theorem and the spin-statistics theorem are proved with great
care in Streater and Wightman, \emph{op.~cit.}}

\problems

\textbf{3.1} \emph{Lorentz group.}  Recall from Eq.~\eqref{eq:3.17} the
Lorentz commutation relations,
\begin{equation}
  [J^{\mu\nu}, J^{\rho\sigma}] = i\bigl( g^{\nu\rho}J^{\mu\sigma}
    - g^{\mu\rho}J^{\nu\sigma} - g^{\nu\sigma}J^{\mu\rho}
    + g^{\mu\sigma}J^{\nu\rho} \bigr).
\end{equation}
\begin{enumerate}[label=(\alph*)]
  \item Define the generators of rotations and boosts as
    \begin{equation}
      J^i = \tfrac{1}{2}\epsilon^{ijk}J^{jk},\qquad
      K^i = J^{0i},
    \end{equation}
    where $i, j, k = 1, 2, 3$.  An infinitesimal Lorentz transformation can
    then be written
    \begin{equation}
      \Phi \to (1 - i\bm{\theta}\cdot\mathbf{J} - i\bm{\beta}\cdot\mathbf{K})\Phi .
    \end{equation}
    Write the commutation relations of these vector operators explicitly.
    (For example, $[J^i, J^j] = i\epsilon^{ijk}J^k$, $[K^i, J^j] =
    i\epsilon^{ijk}K^k$, $[K^i, K^j] = -i\epsilon^{ijk}J^k$.)  Show that the
    combinations
    \begin{equation}
      \mathbf{J}^{+} = \tfrac{1}{2}(\mathbf{J} + i\mathbf{K})
      \qquad\text{and}\qquad
      \mathbf{J}^{-} = \tfrac{1}{2}(\mathbf{J} - i\mathbf{K})
    \end{equation}
    commute with one another and separately satisfy the commutation relations
    of angular momentum.
  \item The finite-dimensional representations of the rotation group
    correspond precisely to the allowed values for angular momentum: integers
    or half-integers.  The result of part (a) implies that all
    finite-dimensional representations of the Lorentz group correspond to
    pairs of integers or half integers, $(j_+, j_-)$, corresponding to pairs
    of representations of the rotation group.  Using the fact that
    $J = \sigma/2$ in the spin-$1/2$ representation of angular momentum,
    write explicitly the transformation laws of the 2-component objects
    transforming according to the $(\tfrac{1}{2}, 0)$ and
    $(0, \tfrac{1}{2})$ representations of the Lorentz group.  Show that
    these correspond precisely to the transformations of $\psi_L$ and
    $\psi_R$ given in \eqref{eq:3.39}.
  \item The identity $\sigma^T = -\sigma^2\sigma\sigma^2$ allows us to
    rewrite the $\psi_L$ transformation in the unitarily equivalent form
    \begin{equation}
      \psi_L \to \psi_L'\Bigl( 1 + i\bm{\theta}\cdot\tfrac{\bm{\sigma}}{2}
        + \bm{\beta}\cdot\tfrac{\bm{\sigma}}{2} \Bigr),
    \end{equation}
    where $\psi_L' = \psi_L^{\dagger}\sigma^2$.  Using this law, we can
    represent the object that transforms as $(\tfrac{1}{2}, \tfrac{1}{2})$ as
    a $2\times2$ matrix that has the $\psi_R$ transformation law on the left
    and, simultaneously, the transposed $\psi_L$ transformation on the right.
    Parametrize this matrix as
    \begin{equation}
      V = \begin{pmatrix}
        V^0 + V^3 & V^1 - iV^2\\
        V^1 + iV^2 & V^0 - V^3
      \end{pmatrix}.
    \end{equation}
    Show that the object $V^\mu$ transforms as a 4-vector.
\end{enumerate}

\textbf{3.2} Derive the Gordon identity,
\begin{equation}
  \bar{u}(p')\gamma^\mu u(p) = \bar{u}(p')\Bigl(
    \frac{p'^{\mu} + p^{\mu}}{2m} + \frac{i\sigma^{\mu\nu}q_\nu}{2m}
    \Bigr)u(p),
\end{equation}
where $q = (p' - p)$.  We will put this formula to use in Chapter~6.

\textbf{3.3} \emph{Spinor products.}  (This problem, together with
Problems~5.3 and~5.6, introduces an efficient computational method for
processes involving massless particles.)  Let $k_0$, $k_1$ be fixed 4-vectors
satisfying $k_0^2 = 0$, $k_1^2 = -1$, $k_0\cdot k_1 = 0$.  Define basic
spinors in the following way: Let $u_{L0}$ be the left-handed spinor for a
fermion with momentum $k_0$.  Let $u_{R0} = k_{1L}u_{L0}$.  Then, for any
$p$ such that $p$ is lightlike ($p^2 = 0$), define
\begin{equation}
  u_L(p) = \frac{1}{\sqrt{2p\cdot k_0}}\,p_{R}u_{L0},\qquad
  \text{and}\qquad
  u_R(p) = \frac{1}{\sqrt{2p\cdot k_0}}\,p_{L}u_{R0}.
\end{equation}
This set of conventions defines the phases of spinors unambiguously (except
when $p$ is parallel to $k_0$).
\begin{enumerate}[label=(\alph*)]
  \item Show that $\slashed{k}_0 u_{R0} = 0$.  Show that, for any lightlike
    $p$, $\slashed{p}u_L(p) = \slashed{p}u_R(p) = 0$.
  \item For the choices $k_0 = (E, 0, 0, -E)$, $k_1 = (0, 1, 0, 0)$,
    construct $u_{R0}$, $u_L(p)$, and $u_R(p)$ explicitly.
  \item Define the spinor products $s(p_1, p_2)$ and $t(p_1, p_2)$, for
    $p_1, p_2$ lightlike, by
    \begin{equation}
      s(p_1, p_2) = \bar{u}_R(p_1)u_L(p_2),\qquad
      t(p_1, p_2) = \bar{u}_L(p_1)u_R(p_2) .
    \end{equation}
    Using the explicit forms for the $u$'s given in part (b), compute the
    spinor products explicitly and show that
    $t(p_1, p_2) = (s(p_2, p_1))^{*}$ and $s(p_1, p_2) = -s(p_2, p_1)$.
    In addition, show that
    \begin{equation}
      |s(p_1, p_2)|^2 = 2p_1\cdot p_2 .
    \end{equation}
    Thus the spinor products are the square roots of 4-vector dot products.
\end{enumerate}

\textbf{3.4} \emph{Majorana fermions.}  Recall from Eq.~\eqref{eq:3.46} that
one can write a relativistic equation for a massless 2-component fermion
field that transforms as the upper two components of a Dirac spinor
$(\tfrac{1}{2}, 0)$.  Call such a 2-component field $\chi_\alpha(x)$,
$\alpha = 1, 2$.
\begin{enumerate}[label=(\alph*)]
  \item Show that it is possible to write an equation for $\chi(x)$ as a
    massive field in the following way:
    \begin{equation}
      i\bar{\sigma}\cdot\partial\,\chi - im\sigma^2\chi^{*} = 0 .
    \end{equation}
    That is, show, first, that this equation is relativistically invariant
    and, second, that it implies the Klein--Gordon equation,
    $(\partial^2 + m^2)\chi = 0$.  This form of the fermion mass is called a
    \emph{Majorana mass term}.
  \item Does the Majorana equation follow from a Lagrangian?  The mass term
    would seem to be the variation of $(\sigma^2)_{ab}\chi_a\chi_b$; however,
    since $\sigma^2$ is antisymmetric, this expression would vanish if
    $\chi(x)$ were an ordinary c-number field.  When we go to quantum field
    theory, we know that $\chi(x)$ will become an anticommuting quantum
    field.  Therefore, it makes sense to develop its classical theory by
    considering $\chi(x)$ as a classical \emph{anticommuting} field, that is,
    as a field that takes as values Grassmann numbers which satisfy
    \begin{equation}
      \alpha\beta = -\beta\alpha
    \end{equation}
    for any $\alpha, \beta$.  Note that this relation implies that
    $\alpha^2 = 0$.  A Grassmann field $\zeta(x)$ can be expanded in a basis
    of functions as
    \begin{equation}
      \zeta(x) = \sum_n a_n\phi_n(x),
    \end{equation}
    where the $\phi_n(x)$ are orthogonal c-number functions and the $a_n$ are
    a set of independent Grassmann numbers.  Define the complex conjugate of
    a product of Grassmann numbers to reverse the order:
    \begin{equation}
      (\alpha\beta)^{*} = \beta^{*}\alpha^{*} = -\alpha^{*}\beta^{*} .
    \end{equation}
    This rule imitates the Hermitian conjugation of quantum fields.  Show
    that the classical action,
    \begin{equation}
      S = \int d^4x\ \Bigl[ \chi^{\dagger}i\bar{\sigma}\cdot\partial\chi
        + \tfrac{im}{2}\bigl( \chi^T\sigma^2\chi - \chi^{\dagger}\sigma^2
          \chi^{*} \bigr) \Bigr]
    \end{equation}
    (where $\chi^{\dagger} = (\chi^{*})^T$) is real ($S^{*} = S$), and that
    varying this $S$ with respect to $\chi$ and $\chi^{*}$ yields the
    Majorana equation.
  \item Let us write a 4-component Dirac field as
    \begin{equation}
      \psi(x) = \begin{pmatrix}\chi_1(x)\\ \chi_2(x)\end{pmatrix},
    \end{equation}
    and recall that the lower components of $\psi$ transform in a way
    equivalent by a unitary transformation to the complex conjugate of the
    representation $(\tfrac{1}{2}, 0)$.  In this way, we can rewrite the
    4-component Dirac field in terms of two 2-component spinors:
    \begin{equation}
      \psi_L(x) = \chi_1(x),\qquad
      \psi_R(x) = i\sigma^2\chi_2^{*}(x) .
    \end{equation}
    Rewrite the Dirac Lagrangian in terms of $\chi_1$ and $\chi_2$ and note
    the form of the mass term.
  \item Show that the action of part (c) has a global symmetry.  Compute the
    divergences of the currents
    \begin{equation}
      J^\mu = \chi^{\dagger}\bar{\sigma}^\mu\chi,\qquad
      J^\mu = \chi_1^{\dagger}\bar{\sigma}^\mu\chi_1
        - \chi_2^{\dagger}\bar{\sigma}^\mu\chi_2,
    \end{equation}
    for the theories of parts (b) and (c), respectively, and relate your
    results to the symmetries of these theories.  Construct a theory of $N$
    free massive 2-component fermion fields with $O(N)$ symmetry (that is,
    the symmetry of rotations in an $N$-dimensional space).
  \item Quantize the Majorana theory of parts (a) and (b).  That is, promote
    $\chi(x)$ to a quantum field satisfying the canonical anticommutation
    relation
    \begin{equation}
      \{\chi_\alpha(\mathbf{x}), \chi_\beta^{\dagger}(\mathbf{y})\}
        = \delta_{\alpha\beta}\delta^{(3)}(\mathbf{x}-\mathbf{y}),
    \end{equation}
    construct a Hermitian Hamiltonian, and find a representation of the
    canonical commutation relations that diagonalizes the Hamiltonian in
    terms of a set of creation and annihilation operators.  (\emph{Hint}:
    Compare $\chi(x)$ to the top two components of the quantized Dirac
    field.)
\end{enumerate}

\textbf{3.5} \emph{Supersymmetry.}  It is possible to write field theories
with continuous symmetries linking fermions and bosons; such transformations
are called supersymmetries.
\begin{enumerate}[label=(\alph*)]
  \item The simplest example of a supersymmetric field theory is the theory
    of a free complex boson and a free Weyl fermion, written in the form
    \begin{equation}
      \mathcal{L} = \partial_\mu\phi^{*}\partial^\mu\phi
        + \chi^{\dagger}i\bar{\sigma}\cdot\partial\chi + F^{*}F .
    \end{equation}
    Here $F$ is an auxiliary complex scalar field whose field equation is
    $F = 0$.  Show that this Lagrangian is invariant (up to a total
    divergence) under the infinitesimal transformation
    \begin{align}
      \delta\phi &= -i\epsilon^T\sigma^2\chi,\\
      \delta\chi &= \epsilon F + \sigma\cdot\partial\phi\,\sigma^2\epsilon^{*},\\
      \delta F &= -i\epsilon^{\dagger}\bar{\sigma}\cdot\partial\chi,
    \end{align}
    where the parameter $\epsilon_\alpha$ is a 2-component spinor of
    Grassmann numbers.
  \item Show that the term
    \begin{equation}
      \Delta\mathcal{L} = \bigl[ m\phi F + \tfrac{i}{2}m\chi^T\sigma^2\chi \bigr]
        + (\text{complex conjugate})
    \end{equation}
    is also left invariant by the transformation given in part (a).
    Eliminate $F$ from the complete Lagrangian $\mathcal{L} +
    \Delta\mathcal{L}$ by solving its field equation, and show that the
    fermion and boson fields $\phi$ and $\chi$ are given the same mass.
  \item It is possible to write supersymmetric nonlinear field equations by
    adding cubic and higher-order terms to the Lagrangian.  Show that the
    following rather general field theory, containing the fields
    $(\phi_i, \chi_i)$, $i = 1, \ldots, n$, is supersymmetric:
    \begin{equation}
      \mathcal{L} = \partial_\mu\phi_i^{*}\partial^\mu\phi_i
        + \chi_i^{\dagger}i\bar{\sigma}\cdot\partial\chi_i + F_i^{*}F_i
        + \Bigl[ W_i F_i - \tfrac{1}{2}W_{ij}\chi_i\chi_j + \text{c.c.}
          \Bigr],
    \end{equation}
    where $W[\phi]$ is an arbitrary function of the $\phi_i$, called the
    \emph{superpotential}, and subscripts denote differentiation with respect
    to the $\phi_i$.  For the simple case $n = 1$ and $W = g\phi^3/3$, write
    out the field equations for $\phi$ and $\chi$ (after elimination of
    $F$).
\end{enumerate}

\textbf{3.6} \emph{Fierz transformations.}  Let $u_i$, $i = 1, \ldots, 4$, be
four 4-component Dirac spinors.  In the text, we proved the Fierz
rearrangement formulae \eqref{eq:3.83} and \eqref{eq:3.84}.  The first of
these formulae can be written in 4-component notation as
\begin{equation}
  (\bar{u}_1\gamma^\mu P_L u_2)(\bar{u}_3\gamma_\mu P_L u_4)
    = -(\bar{u}_1\gamma^\mu P_L u_4)(\bar{u}_3\gamma_\mu P_L u_2)
      + \cdots .
\end{equation}
In fact, there are similar rearrangement formulae for any product
\begin{equation}
  (\bar{u}_1\Gamma^A u_2)(\bar{u}_3\Gamma^B u_4),
\end{equation}
where $\Gamma^A, \Gamma^B$ are any of the 16 combinations of Dirac matrices
listed in Section~3.4.
\begin{enumerate}[label=(\alph*)]
  \item To begin, normalize the 16 matrices $\Gamma^A$ to the convention
    \begin{equation}
      \mathrm{tr}[\Gamma^A\Gamma^B] = 4\delta^{AB} .
    \end{equation}
    This gives $\Gamma^A = \{1, \gamma^0, i\gamma^1, \ldots\}$; write all 16
    elements of this set.
  \item Write the general Fierz identity as an equation
    \begin{equation}
      (\bar{u}_1\Gamma^A u_2)(\bar{u}_3\Gamma^B u_4) =
        \sum_{C,D} C^{ABCD}(\bar{u}_1\Gamma^C u_4)(\bar{u}_3\Gamma^D u_2),
    \end{equation}
    with unknown coefficients $C^{ABCD}$.  Using the completeness of the 16
    $\Gamma^A$ matrices, show that
    \begin{equation}
      C^{ABCD} = \tfrac{1}{16}\mathrm{tr}[\Gamma^C\Gamma^A\Gamma^D\Gamma^B] .
    \end{equation}
  \item Work out explicitly the Fierz transformation laws for the products
    $(\bar{u}_1u_2)(\bar{u}_3u_4)$ and
    $(\bar{u}_1\gamma^\mu u_2)(\bar{u}_3\gamma_\mu u_4)$.
\end{enumerate}

\textbf{3.7} This problem concerns the discrete symmetries $P$, $C$, and $T$.
\begin{enumerate}[label=(\alph*)]
  \item Compute the transformation properties under $P$, $C$, and $T$ of the
    antisymmetric tensor fermion bilinears,
    \begin{equation}
      \bar{\psi}\sigma^{\mu\nu}\psi,\qquad
      \text{with } \sigma^{\mu\nu} = \tfrac{i}{2}[\gamma^\mu, \gamma^\nu] .
    \end{equation}
    This completes the table of the transformation properties of bilinears at
    the end of the chapter.
  \item Let $\phi(x)$ be a complex-valued Klein--Gordon field, such as we
    considered in Problem~2.2.  Find unitary operators $P$, $C$ and an
    antiunitary operator $T$ (all defined in terms of their action on the
    annihilation operators $a_{\mathbf{p}}$ and $b_{\mathbf{p}}$ for the
    Klein--Gordon particles and antiparticles) that give the following
    transformations of the Klein--Gordon field:
    \begin{equation}
      P\phi(t,\mathbf{x})P = \phi(t,-\mathbf{x});\qquad
      T\phi(t,\mathbf{x})T = \phi(-t,\mathbf{x});\qquad
      C\phi(t,\mathbf{x})C = \phi^{*}(t,\mathbf{x}) .
    \end{equation}
    Find the transformation properties of the components of the current
    \begin{equation}
      j^\mu = i\bigl( \phi^{*}\partial^\mu\phi - \phi\partial^\mu\phi^{*} \bigr)
    \end{equation}
    under $P$, $C$, and $T$.
  \item Show that any Hermitian Lorentz-scalar local operator built from
    $\psi(x)$, $\phi(x)$, and their conjugates has $CPT = +1$.
\end{enumerate}

\textbf{3.8} \emph{Bound states.}  Two spin-$1/2$ particles can combine to a
state of total spin either 0 or 1.  The wavefunctions for these states are
odd and even, respectively, under the interchange of the two spins.
\begin{enumerate}[label=(\alph*)]
  \item Use this information to compute the quantum numbers under $P$ and $C$
    of all electron-positron bound states with $S$, $P$, or $D$ wavefunctions.
  \item Since the electron-photon coupling is given by the Hamiltonian
    \begin{equation}
      H_I = -e\int d^3x\ \bar{\psi}\gamma^\mu\psi\,A_\mu,
    \end{equation}
    where $j^\mu$ is the electric current, electrodynamics is invariant to
    $P$ and $C$ if the components of the vector potential have the same $P$
    and $C$ parity as the corresponding components of $j^\mu$.  Show that
    this implies the following surprising fact: The spin-0 ground state of
    positronium can decay to 2 photons, but the spin-1 ground state must
    decay to 3 photons.  Find the selection rules for the annihilation of
    higher positronium states, and for 1-photon transitions between
    positronium levels.
\end{enumerate}
